
\section{Cohomological descent to a shifted Poisson vertex algebra}
\label{sec:PVA_descent}
\label{sec:Ainfty-to-PVA}

\paragraph{Discipline: classical $\lambda$-Jacobi versus quantum lift.}
Classical $\lambda$-Jacobi: \ClaimStatusProvedHere{} (PVA descent at the
classical level closes unconditionally on the proved scope of
Definition~\ref{def:log-SC-algebra}). All-loop quantum lift:
\ClaimStatusConditional{} on the four endpoint hypotheses
$\hypKZSDR$ (Khan--Zeng analytic spectral-decomposition representation
\cite{KZ25}), $\hypStokes$ (Stokes-data choice on the Fulton--MacPherson
compactification), $\hypReflWts$ (reflected weights / wall-crossing
under analytic continuation), and $\hypTLift$ (chain-level lift of
$T = [Q_{\mathrm{tot}}, G]$ to the all-loop quantum object;
Costello \cite{Costello2011Renorm}; Costello--Gwilliam \cite{CG17};
Lurie HA chap.~5 \cite{LurieHA}). The classical-to-quantum bridge is
the all-loop PVA--quantum lifting problem. Theorem statements that depend
on the all-loop quantum lift carry the four $\hyp$ tags; classical
statements remain unconditional.

The $A_\infty$ chiral algebra $(\cA, Q, \{m_k\})$ is a chain-level
object: the operations $m_k$ act on cochains, not on cohomology
classes. Passing to cohomology $H^\bullet(\cA, Q)$ isolates the
binary classical shadow: the surviving binary operation~$m_2$
descends to a $(-1)$-shifted Poisson vertex algebra, with a
commutative product (from the regular part of~$m_2$) and a
skew-symmetric $\lambda$-bracket (from the singular part)
satisfying Jacobi and Leibniz. Vanishing of the higher operations
$m_{k\ge3}$ is a separate formality statement, proved below only
under a compatible topological nullhomotopy hypothesis. This
descent is the motivic projection of
Remark~\ref{rem:pva-depth-decomposition}: it retains the arithmetic
content of the bar complex and forgets the higher homotopy data not
seen by the PVA.

\begin{remark}[PVA descent as the chain-level avatar of the bar complex
at the quasi-classical limit]
\label{rem:pva-quasi-classical-bar-avatar}
The PVA descent of Theorem~\ref{thm:cohomology-PVA-main} is the
\emph{chain-level avatar of the bar complex at the quasi-classical
limit} $\hbar \to 0$ of the BV quantization (cf.\ Remark~\ref{rem:bv-hbar-convention}
of Section~\ref{sec:BV_construction}). Concretely: the chiral bar complex
$\barB^{\mathrm{ch}}(\cA) = T^c(s^{-1}\bar\cA)$ carries a Maurer--Cartan
element $\Theta_\cA$ satisfying $D_\cA\Theta_\cA + \tfrac12[\Theta_\cA,\Theta_\cA] = 0$
(universal-holography master, Theorem~\ref{thm:universal-holography-master});
the binary operation $m_2$ is the leading component of $\Theta_\cA$ at bar
degree~$2$, and the $(-1)$-shifted PVA structure $(\mu, \{-_\lambda-\})$
on $H^\bullet(\cA,Q)$ is the projection of $\Theta_\cA$ to the classical
slice. In the language of the Volume~I universal-conductor theorem
($K = -c_{\mathrm{ghost}}(\mathrm{BRST})$), the PVA bracket retains the
\emph{ghost-charge--linear} component of $\Theta_\cA$: the $\lambda$-bracket
encodes the BRST one-cocycle generated by $\hbar\Delta_{\mathrm{BV}}$ acting
on $S_{\mathrm{eff}}$ at first order in $\hbar$, while the higher Massey
products $m_{k \ge 3}$ encode higher-order quantum corrections outside
the binary PVA shadow unless a compatible nullhomotopy kills them. The
low-order lambda-bracket coefficients
(verified via Khan--Zeng \cite{KZ25}, Proposition~\ref{prop:abelian-study},
and the Heisenberg / current--algebra checks of \S\ref{subsec:Higgs-Coulomb})
are therefore the chain-level shadows of the universal holography master
theorem at the level of binary operations.
\end{remark}

\begin{remark}[$W(p)$ triplet: scope qualification]
\label{rem:pva-wp-tempering-scope}
The PVA descent of Theorem~\ref{thm:cohomology-PVA-main} applies
unconditionally to logarithmic $\SCchtop$-algebras whose closed-colour
weight forms have at most \emph{logarithmic} singularities
(Definition~\ref{def:log-SC-algebra}), covering the standard
non-logarithmic landscape (Heisenberg / class~G, affine Kac--Moody at
non-critical level / class~L, $\beta\gamma$ / class~C, Virasoro and
$W_N$ at generic central charge / class~M). For the logarithmic
$W(p)$ triplet algebra, four strong generators are identified ($T$
together with the $\mathfrak{sl}_2$ triplet $W^a$,
$a \in \{-1, 0, 1\}$); whether these constitute a \emph{free} strong
generation system compatible with the Khan--Zeng PVA framework
\cite{KZ25} remains open. Consequently, the $W(p)$ triplet tempering
requires the Adamovi\'c--Milas character-amplitude bound: the
chain-level PVA descent has the sharp harmonic-decoupling obstruction
\(\mathfrak h_{g,r}^{\mathsf C}\)
of Definition~\ref{def:class-c-harmonic-decoupling-obstruction} at the
quartic-pole stratum present in $W(p)$. The $\hbar \to 0$
classical limit and the
descended $\lambda$-bracket exist as formal series in all cases; the
open issue is whether the resulting series converges in the
weight-completed category for $W(p)$, equivalently whether the
classes \(\mathfrak h_{g,r}^{\mathsf C}\) vanish and lift through the
logarithmic stratum. For all
non-logarithmic families in the standard-landscape universal-holography
scope, the descent is unconditional on the proved scope.
\end{remark}

\begin{remark}[Recognition theorem; PVA as detection invariant]
\label{rem:pva-recognition-detection}
The recognition theorem (Theorem~\ref{thm:recognition-SC}, \S\ref{subsec:from-prefact},
Theorem~\ref{thm:Obs-is-SC}) identifies which logarithmic $\SCchtop$-algebra
arises from a given chiral algebra; the PVA descent of
Theorem~\ref{thm:cohomology-PVA-main} provides the \emph{detection invariant}
on the cohomology side: the $(-1)$-shifted PVA $(H^\bullet(\cA,Q), \mu, \{-_\lambda-\})$
recovers the classical limit of the chiral algebra together with its central
charge data ($\kappa$ via the Sugawara stress tensor, $c$ via the Virasoro
sub-algebra). When two boundary algebras $A_\partial$ and $A'_\partial$ produce
quasi-isomorphic $\SCchtop$-algebras under the recognition theorem, their
PVA descents agree. The converse is only a detection statement:
different PVA shadows prove inequivalence, while equal PVA shadows
leave the ordered $R$-matrix, higher $A_\infty$ operations, and
bulk--boundary action to be checked. This is the chain-level analogue of the universal
holography mirror principle: chiral input $\leftrightarrow$ bulk
$\SCchtop$-output $\leftrightarrow$ PVA classical limit.
\end{remark}

Three points of the construction require emphasis.

\begin{enumerate}[label=\textup{(\arabic*)}]
\item We work with the \emph{unsymmetrized} regular and singular parts of the binary kernel.
 Graded commutativity and skew-symmetry are then \emph{theorems}, not definitions.
\item We separate the two formal variables that play different roles. The collision/OPE
 variable is denoted by $\zeta$, while the standard polynomial PVA variable is denoted
 by $\lambda$. The two are related by the usual Borel transform of singular coefficients.
\item The Jacobi and Leibniz identities are derived from the \emph{three} pairwise
 boundary divisors $D_{\{1,2\}},D_{\{1,3\}},D_{\{2,3\}}\subset \partial\FM_3(\C)$.
 The non-consecutive face $D_{\{1,3\}}$ is not discarded; it is precisely the geometric
 source of the third Jacobi/Leibniz term.
\end{enumerate}

\begin{maintheorem}[Cohomological descent to a shifted PVA; \ClaimStatusProvedHere]
\label{thm:cohomology-PVA-main}%
\label{thm:cohomology_PVA}%
\label{thm:pva-coisson-bridge}%
Let $(\A,Q,\partial,\mathbf 1,\{m_k\}_{k\ge1})$ be a
logarithmic\/ $\SCchtop$-algebra
\textup{(}Definition~\textup{\ref{def:log-SC-algebra})}.
Write the binary kernel in the collision variable $\zeta$ as
\[
m_2(a,b;\zeta)=m_2^{\mathrm{reg}}(a,b;\zeta)+m_2^{\mathrm{sing}}(a,b;\zeta),
\qquad
m_2^{\mathrm{sing}}(a,b;\zeta)=\sum_{n\ge0}(a_{(n)}b)\,\zeta^{-n-1}.
\]
Define the unsuspended regular product and the singular
$\lambda$-bracket by
\[
\mu(a,b):=\mu_2^{\mathrm{reg}}(a,b;0),
\qquad
\{a_\lambda b\}:=\sum_{n\ge0}\frac{\lambda^n}{n!}\,a_{(n)}b.
\]
Here $\mu_2^{\mathrm{reg}}$ is the ordinary regular product whose
suspended bar component is $m_2^{\mathrm{reg}}$.
Then:

\begin{enumerate}[label=\textup{(\roman*)}]
\item The operations $\mu$ and $\{-{}_\lambda-\}$ descend to $H^\bullet(\A,Q)$.
\item The descended product is unital, associative, and graded commutative; the descended
 $\lambda$-bracket has cohomological degree $-1$ and satisfies sesquilinearity,
 ordinary vertex-algebra skew-symmetry, Jacobi, and the Leibniz rule
 with the shifted Koszul signs. Hence
 $H^\bullet(\A,Q)$ is a $(-1)$-shifted Poisson vertex algebra.
\end{enumerate}
\end{maintheorem}

The rest of the section proves these two parts and then records the
separate higher-vanishing criterion.

\begin{remark}[Depth decomposition and PVA formality]
\label{rem:pva-depth-decomposition}
The PVA descent is not a formality theorem for the whole
$A_\infty$ algebra. In the language of Vol~I
(via the depth decomposition theorem), it is the projection to the
binary classical shadow: the product and the $\lambda$-bracket
survive, while the higher Massey products measure algebraic depth
outside the PVA. The chiral algebra may have
$d_{\mathrm{alg}} = \infty$ (e.g., Virasoro, where
$T_{(1)}T = 2T$ drives an infinite shadow cascade), but the PVA
shadow retains only the \emph{arithmetic} depth
$d_{\mathrm{arith}}$ (Hecke eigenforms, $L$-function content).
On an additional formal or topologically nullhomotopic locus the
higher operations become $Q$-exact, but that is the separate
criterion of Proposition~\ref{prop:m3_vanish}. The PVA descent is
therefore a \emph{motivic projection}: it forgets the mixed-Tate
(homotopy-theoretic) part and preserves
the non-Tate (arithmetic) part of the depth decomposition, namely
the binary product $m_2$ and its residues.
In the shadow metric formulation
(of Volume~I),
the PVA shadow is the
$\hbar \to 0$ limit of $Q_L(t)$: the critical discriminant
$\Delta = 8\kappa S_4$ vanishes on the classical slice,
recovering the Gaussian envelope
$Q_L^{\mathrm{cl}}(t) = (2\kappa + 3\alpha t)^2$.
\end{remark}

\begin{remark}[PVA descent as the classical projection of a two-stage tower]
\label{rem:pva-two-stage-projection}
\index{PVA descent!two-stage projection}
\index{averaging map!PVA as further projection}
The Poisson vertex algebra $H^\bullet(\cA, Q)$ arises by a
two-stage projection from the ordered (primitive) bar data. The
first stage is the averaging map
$\mathrm{av}\colon T^c(s^{-1}\bar\cA) \to \Sym^c(s^{-1}\bar\cA)$,
which passes from the ordered bar complex $\barB^{\mathrm{ord}}(\cA)$
(the $E_1$ coalgebra, retaining full $R$-matrix and spectral data)
to the symmetric bar complex $\barB^{\Sigma}(\cA)$ (the modular
coalgebra, retaining the scalar invariants $\kappa$, cubic shadow, etc.).
The second stage is the passage to cohomology
$H^\bullet(\cA, Q)$: the binary operation~$m_2$ descends to the
commutative product and $\lambda$-bracket of the PVA, while the
higher $A_\infty$ operations $m_{k \ge 3}$ are forgotten by the
binary classical shadow unless the additional nullhomotopy criterion
of Proposition~\ref{prop:m3_vanish} applies.
Thus the full chain of projections is
\[
 T^c \;\xrightarrow{\;\mathrm{av}\;}
 \Sym^c \;\xrightarrow{\;H^\bullet\;}
 \mathrm{PVA}.
\]
The first arrow forgets the $R$-matrix (ordered-to-unordered); the
second forgets homotopy (chain-level to cohomology). The
PVA retains the arithmetic content of the bar complex
(binary residues, Hecke eigenforms) while discarding both
the spectral data ($R$-matrix, Yangian coherences) and the
homotopy data (higher $A_\infty$ operations, shadow obstruction
tower).
\end{remark}

\begin{remark}[Unsymmetrized vs.\ symmetrized product]
\label{rem:unsymmetrized-product}
The product $\mu(a,b) := \mu_2^{\mathrm{reg}}(a,b;0)$ defined above is
\emph{unsymmetrized}: it uses the ordinary regular product
$\mu_2^{\mathrm{reg}}(a,b;0)$ directly. Its suspended bar component is
$m_2^{\mathrm{reg}}$.  Definition~\ref{def:product} in
\S\ref{subsec:product-bracket} records the same cohomology class in
suspended notation; after desuspension it is represented by
$[a]\cdot[b] := \tfrac{1}{2}\bigl(\mu_2^{\mathrm{reg}}(a,b;0) +
(-1)^{|a||b|}\mu_2^{\mathrm{reg}}(b,a;0)\bigr)$.
The two agree on $H^\bullet(\A,Q)$: the exchange cylinder
(Proposition~\ref{prop:product-commutative}) proves graded commutativity of
$\mu$ on cohomology, so the symmetrization is redundant \emph{a~posteriori}.
The unsymmetrized form simplifies the chain-level
Stokes arguments; the symmetrized form in \S\ref{subsec:product-bracket}
makes skew-symmetry manifest at the cochain level.
\end{remark}

\begin{remark}[Symplectic origin of the shifted PVA]
\label{rem:symplectic-origin-PVA}
The $(-1)$-shifted PVA on $H^\bullet(\A,Q)$ is
related to a $(-1)$-shifted symplectic structure (PTVV), via two steps.

\smallskip\noindent\emph{Step~1: from $\mathfrak{S}_b$ to the bar complex.}
The self-intersection
$\mathfrak{S}_b = \mathscr{L}_b \times_{\mathscr{M}} \mathscr{L}_b$ of the
Lagrangian embedding $\mathscr{L}_b \hookrightarrow \mathscr{M}$ carries a
canonical $(-1)$-shifted symplectic form (PTVV).
By Theorem~\ref{thm:bar-is-self-intersection}, the bar complex
$\barB(\A)$ provides the formal bar/Koszul model for the derived
functions on $\mathfrak{S}_b$ near the boundary locus, so the
$(-1)$-shifted symplectic form induces a $(-1)$-shifted Poisson
bracket on $\barB(\A)$.

\smallskip\noindent\emph{Step~2: descent to cohomology.}
The bracket on $\barB(\A)$ descends along the bar-to-boundary
comparison map
\[
\chi\colon H^\bullet(\barB(\A))\longrightarrow H^\bullet(\A,Q).
\]
When the Koszul-effectiveness hypotheses make~$\chi$ an isomorphism,
this is the shifted Darboux bracket on boundary cohomology. Without
that hypothesis it is a Poisson bracket on the bar/Koszul shadow and
the direct FM proof below supplies the PVA bracket on
$H^\bullet(\A,Q)$. It is at this stage that the $\lambda$-bracket
$\{{\cdot}_\lambda{\cdot}\}$ appears, via the residue extraction
along the collision divisor $D_{\{1,2\}} \subset \FM_2(\C)$.

\smallskip
The proofs of the PVA axioms below (Jacobi from the three-face
Stokes argument on $\FM_3(\C)$, Leibniz from the mixed
regular--singular three-face sector, skew-symmetry from the exchange
half-monodromy) are direct
computations using the manifold-with-corners structure of
$\FM_k(\C)$. They can be \emph{interpreted} as reflecting
geometric properties of the self-intersection $\mathfrak{S}_b$,
but the proofs themselves do not invoke derived symplectic
geometry; they are independent of it.
\end{remark}

\providecommand{\Steinb}{\mathfrak{S}}

\begin{proposition}[Steinberg Poisson comparison for the PVA bracket;
\ClaimStatusConditional; licensing $\alpha+\beta+\gamma+\varepsilon$
via boundary chart $b$, the comparison map $\chi$, the derived
shifted-symplectic ambient, and $\effKoszul$]
\label{prop:PVA-from-symplectic}
Assume the boundary chart $\A$ is presented by the bar model of the
Steinberg self-intersection
$\Steinb_b = \mathscr{L}_b \times_{\mathscr{M}} \mathscr{L}_b$
as in Theorem~\textup{\ref{thm:steinberg-presentation}} and
Remark~\textup{\ref{rem:steinberg-lagrangian-origin}}. Assume also
that the bar-to-boundary comparison map
\[
\chi\colon H^\bullet(\barB(\A))\longrightarrow H^\bullet(\A,Q)
\]
is an isomorphism on the Koszul-effective comparison locus
\(\effKoszul\). Then the
$(-1)$-shifted Poisson bracket $\{a_\lambda b\}$ on $H^\bullet(\A,Q)$ is
the Poisson bracket associated to the $(-1)$-shifted symplectic structure
$\omega_{-1}$ on $\Steinb_b$.
\end{proposition}

\begin{proof}
The argument has three steps.

\smallskip
\noindent\textbf{Step~1 (Functions on the self-intersection).}
By Theorem~\ref{thm:bar-is-self-intersection}, the bar complex
$\barB(\A)$ supplies the formal bar/Koszul model for the derived
functions on the Lagrangian self-intersection near the boundary
locus.
The Lagrangian embedding
$\mathscr{L}_b \hookrightarrow \mathscr{M}$ into the $(-2)$-shifted
symplectic target endows the derived self-intersection $\Steinb_b$ with a
canonical $(-1)$-shifted symplectic form $\omega_{-1}$
(Pantev--To\"en--Vaqui\'e--Vezzosi~\cite{PTVV13}).

\smallskip
\noindent\textbf{Step~2 (From symplectic form to Poisson bracket).}
The $(-1)$-shifted symplectic form $\omega_{-1}$ on $\Steinb_b$ induces,
by the shifted Darboux formalism, a $(-1)$-shifted Poisson bracket
$\{-,-\}_{\omega}$ on $\mathcal{O}(\Steinb_b) \simeq \barB(\A)$.
A remark on structures: $\barB(\A)$ carries a
\emph{coalgebra} structure (the deconcatenation coproduct from the
bar construction), but $\mathcal{O}(\Steinb_b)$ is a commutative
dg-\emph{algebra} (the algebra of derived functions on the
self-intersection). The identification
$\barB(\A) \simeq \mathcal{O}(\Steinb_b)$ endows $\barB(\A)$ with
a commutative product \emph{in addition to} its coalgebra structure;
the Poisson bracket is a binary operation on this commutative
algebra structure, not on the coalgebra. (The two structures are
compatible: they assemble into a bialgebra in the derived sense,
reflecting the Hopf-algebraic nature of the groupoid algebra
$\mathcal{O}(\Steinb_b)$.)
On cohomology the bracket descends through the comparison map~$\chi$;
if $\chi$ is an isomorphism, this gives a well-defined
$(-1)$-shifted Poisson structure on $H^\bullet(\A,Q)$. Without this
effectiveness, the Steinberg bracket remains a Poisson bracket on the
bar/Koszul shadow, while the direct FM proof below supplies the PVA
bracket on \(H^\bullet(\A,Q)\).

\smallskip
\noindent\textbf{Step~3 (Identification with the $\lambda$-bracket).}
The $\lambda$-bracket
\[
\{a_\lambda b\}
\;=\;
\Res_{\zeta=0}\bigl[m_2^{\mathrm{sing}}(a,b;\zeta)\cdot e^{\lambda\zeta}\bigr]
\]
extracts the residue of the binary singular kernel along the collision
divisor $D_{\{1,2\}} \subset \FM_2(\C)$. Under the identification
$\barB(\A) \simeq \mathcal{O}(\Steinb_b)$, this collision divisor is
the conormal direction to the diagonal
$\Delta \colon \mathscr{L}_b \hookrightarrow
\mathscr{L}_b \times_{\mathscr{M}} \mathscr{L}_b$:
the residue along $D_{\{1,2\}}$ extracts the conormal component
of the Poisson bivector, that is, the component in
$N^*_{\Delta/\Steinb_b}$ transverse to the diagonal
$\Delta \colon \mathscr{L}_b \hookrightarrow \Steinb_b$.
The Lagrangian condition identifies this conormal bundle with the
tangent bundle $T\mathscr{L}_b$, so the residue is a bivector on
$\mathscr{L}_b$, which is precisely the $(-1)$-shifted Poisson
bivector $\pi_\omega$ dual to $\omega_{-1}$.
The exponential $e^{\lambda\zeta}$ is the Borel kernel that converts
the Laurent residue into the polynomial $\lambda$-bracket.
Hence $\{a_\lambda b\}$ on $H^\bullet(\A,Q)$ coincides with the
descended Poisson bracket $\{-,-\}_{\omega}$.
\end{proof}

\subsection{Binary descent from the degree-$2$ relation}

\begin{proposition}[$A_\infty$ degree-$2$ compatibility;
\ClaimStatusProvedHere; licensing $\gamma$ via the logarithmic
chain/pro ambient $\hypAmbientPro$]
\label{prop:m1_m2}
The $n=2$ specialization of the $A_\infty$ identity \eqref{eq:ainfty-relation-raw} is
\begin{equation}
\label{eq:repaired-n2}
Q(m_2(a,b))+m_2(Qa,b)+(-1)^{\degree{a}}m_2(a,Qb)=0.
\end{equation}
Equivalently,
\[
Q(m_2(a,b))=-m_2(Qa,b)-(-1)^{\degree{a}}m_2(a,Qb),
\]
so \(Q=m_1\) is compatible with the binary operation in the
desuspended Stasheff convention.
\end{proposition}

\begin{proof}
This is exactly the $(i,j)=(1,2)$ and $(i,j)=(2,1)$ part of
\eqref{eq:ainfty-relation-raw}. The three terms are
\[
m_1(m_2(a,b)),\qquad m_2(m_1(a),b),\qquad
(-1)^{\degree{a}}m_2(a,m_1(b)).
\]
No other term occurs at degree~$2$.
\end{proof}

\begin{lemma}[Descent of the regular and singular binary parts;
\ClaimStatusProvedHere; licensing $\gamma$ via the logarithmic
chain/pro ambient $\hypAmbientPro$]
\label{lem:lambda_descends}
Work in the logarithmic chain/pro ambient $\hypAmbientPro$, with the
regular/singular decomposition of \(m_2\) from
\cref{thm:cohomology-PVA-main} and the polynomial Borel transform of
\cref{def:borel-transform-pva}. If \(a,b\in\A\) are \(Q\)-closed, then
\(\mu(a,b)\) and \(\{a_\lambda b\}\) are \(Q\)-closed, and their
cohomology classes depend only on the classes
\([a],[b]\in H^\bullet(\A,Q)\).
\end{lemma}

\begin{proof}
Since $Q$ acts on coefficients in $\A$ and does not touch the formal variables,
it commutes with the regular/singular projections and with taking the constant
term at $\zeta=0$. Thus \eqref{eq:repaired-n2} implies
\[
Q\mu(a,b)=-\mu(Qa,b)-(-1)^{\degree{a}}\mu(a,Qb),
\]
and, coefficientwise in the singular expansion,
\[
Q(a_{(n)}b)=-(Qa)_{(n)}b-(-1)^{\degree{a}}a_{(n)}(Qb).
\]
Hence $Qa=Qb=0$ implies $Q\mu(a,b)=0$ and $Q(a_{(n)}b)=0$ for all $n\ge0$, so the
polynomial bracket $\{a_\lambda b\}$ is $Q$-closed coefficientwise.

If $a'=a+Qx$, then \eqref{eq:repaired-n2} gives
\[
m_2(a',b)-m_2(a,b)
=
m_2(Qx,b)
=
-Qm_2(x,b)-(-1)^{\degree{x}}m_2(x,Qb).
\]
For $Qb=0$, the difference is $Q$-exact. Applying the regular projection, constant-term
projection, or any coefficient projection of the singular part preserves $Q$-exactness.
The same argument applies in the second variable.
\end{proof}

\begin{proposition}[Sesquilinearity in cohomology;
\ClaimStatusProvedHere; licensing $\gamma$ via the logarithmic
chain/pro ambient $\hypAmbientPro$]
\label{prop:PVA1_proof}
In the polynomial divided-power convention of
\cref{def:borel-transform-pva}, the descended $\lambda$-bracket
satisfies the standard PVA sesquilinearity relations
\begin{align}
\label{eq:repaired-PVA1a}
\{\partial a_\lambda b\} &= -\lambda\,\{a_\lambda b\},\\
\label{eq:repaired-PVA1b}
\{a_\lambda \partial b\} &= (\lambda+\partial)\,\{a_\lambda b\}.
\end{align}
\end{proposition}

\begin{proof}
Write
\[
m_2^{\mathrm{sing}}(a,b;\zeta)
=\sum_{n\ge0}(a_{(n)}b)\,\zeta^{-n-1}.
\]
For \(k=2\), the chain-level sesquilinearity identities
\eqref{eq:sesqui-left}--\eqref{eq:sesqui-right} are equivalent,
coefficientwise, to
\[
(\partial a)_{(0)}b=0,\qquad
(\partial a)_{(n)}b=-n\,a_{(n-1)}b\quad(n\ge1),
\]
and
\[
a_{(n)}(\partial b)=\partial(a_{(n)}b)+n\,a_{(n-1)}b
\quad(n\ge0),
\]
with \(a_{(-1)}b=0\). Applying the polynomial Borel transform gives
\[
\{\partial a_\lambda b\}
=\sum_{n\ge0}(\partial a)_{(n)}b\,\frac{\lambda^n}{n!}
=-\lambda\sum_{m\ge0}a_{(m)}b\,\frac{\lambda^m}{m!}
=-\lambda\{a_\lambda b\},
\]
and
\[
\{a_\lambda\partial b\}
=\sum_{n\ge0}a_{(n)}(\partial b)\,\frac{\lambda^n}{n!}
=\partial\{a_\lambda b\}
+ \lambda\sum_{m\ge0}a_{(m)}b\,\frac{\lambda^m}{m!}
=(\lambda+\partial)\{a_\lambda b\}.
\]
Since \(Q\) commutes with \(\partial\) by
Definition~\ref{def:sesquilinearity} and with the formal variable
\(\lambda\), Lemma~\ref{lem:lambda_descends} carries these identities
to \(H^\bullet(\A,Q)\).
\end{proof}

\subsection{The exchange cylinder and the two-point symmetry theorem}

The $A_\infty$ operations are defined on \emph{ordered} configurations:
the inputs $a_1, \ldots, a_k$ are ordered by $t_1 < \cdots < t_k$
on the topological line. The ordered configuration space
$\Conf_2^{<}(\R)$ is disconnected (its two components correspond
to the two orderings $t_1 < t_2$ and $t_2 < t_1$), so the
topological direction alone cannot exchange two insertions. But
$\Conf_2(\C \times \R)$ \emph{is} connected: two bulk insertions
exchange by winding in the holomorphic plane (a half-turn around the
collision divisor) while simultaneously swapping their topological
coordinates. This exchange cylinder is the geometric mechanism that
produces both graded commutativity of the product and shifted
skew-symmetry of the bracket on cohomology.

\begin{construction}[Explicit exchange path]
\label{const:exchange-cylinder}
Fix $r>0$ and consider two bulk points in $\R_t\times\C_z$:
\[
x_1(0)=(-1,r),\qquad x_2(0)=(1,-r).
\]
Define a path $s\mapsto (x_1(s),x_2(s))\in \Conf_2(\R\times\C)$ by
\begin{equation}
\label{eq:exchange-path}
x_1(s)=\bigl(-\cos(\pi s),\,r\,e^{\pi i s}\bigr),
\qquad
x_2(s)=\bigl(\cos(\pi s),\,-r\,e^{\pi i s}\bigr),
\qquad 0\le s\le1.
\end{equation}
Its separation vector is
\[
x_1(s)-x_2(s)=\bigl(-2\cos(\pi s),\,2r\,e^{\pi i s}\bigr),
\]
which never vanishes: at $s=\tfrac12$ the topological separation vanishes, but the
holomorphic separation is $2ri\neq0$; away from $s=\tfrac12$ the topological separation
is already nonzero. Thus \eqref{eq:exchange-path} is a genuine isotopy in the full
configuration space exchanging the two endpoints.
\end{construction}

\begin{remark}[The exchange cylinder as symplectic rotation]
\label{rem:exchange-symplectic-rotation}
The exchange path~\eqref{eq:exchange-path} traverses the holomorphic
half-boundary arc while simultaneously transposing the topological
$\R$-coordinates. In the
language of the $(-1)$-shifted self-intersection $\mathfrak{S}_b$, this is
the Lagrangian half-monodromy of the correspondence: the holomorphic
half-turn
is the symplectic rotation that connects the two sheets of the
self-intersection, and the topological transposition records the change
of $E_1$-ordering. Commutativity and skew-symmetry on cohomology
(Propositions~\ref{prop:product-commutative} and~\ref{prop:PVA2_proof}) are
consequences of holomorphic half-monodromy around the collision circle
together with topological transposition in the $E_1$ direction.
\end{remark}

\begin{proposition}[Exchange half-monodromy as Lagrangian monodromy;
\ClaimStatusProvedHere; licensing $\alpha+\gamma$ via boundary chart
$b$ and the logarithmic chain/pro ambient $\hypAmbientPro$]
\label{prop:exchange-lagrangian-monodromy}
Work over the boundary chart \(b\) and in the logarithmic chain/pro
ambient \(\hypAmbientPro\).  The exchange path
\textup{(}holomorphic half-monodromy in $\C$ combined with
topological transposition in $\R$\textup{)} is the Lagrangian
sheet-exchange monodromy of the $(-1)$-shifted self-intersection\/
$\mathfrak{S}_b$ around the collision divisor\/
$D_{1,2} \subset \FM_2(\C)$. Specifically:
\begin{enumerate}[label=\textup{(\roman*)}]
\item Its holomorphic separation is
 \(z_1(s)-z_2(s)=2r\,e^{\pi i s}\).  Hence the projection to the
 boundary circle of \(\FM_2(\C)\) traverses the half-boundary arc
 from angle \(0\) to angle \(\pi\), not a loop.  The square of the
 exchange traverses the full boundary circle and represents the
 generator of \(\pi_1(\FM_2(\C))\cong\Z\).  No braid-group \(\Z\) is
 assigned to the full configuration space of two ordered points in
 \(\C\times\R\).
\item The monodromy of the $(-1)$-shifted symplectic form\/
 $\omega_{-1}$ along this holomorphic boundary half-arc is the
 involution $\sigma$ exchanging the two sheets of the
 self-intersection.
\item The skew-symmetry\/
 $\{a_\lambda b\} = -\{b_{-\lambda-\partial} a\}$ is the Borel
 transform of the oriented Stokes boundary identity for the exchange
 cylinder.  The exchange sends \(\zeta=z_1-z_2\) to \(-\zeta\);
 \(d\log(-\zeta)=d\log\zeta\), so the minus sign is the oriented
 boundary sign in the chain homotopy, not a sign of the logarithmic
 form.
\end{enumerate}
\end{proposition}

\begin{proof}
The space $\mathrm{Conf}_2(\R)$ is disconnected, with components
indexed by the order. The exchange path is therefore not a loop in
$\FM_2(\C)\times\mathrm{Conf}_2(\R)$. From
\eqref{eq:exchange-path},
\[
 z_1(s)-z_2(s)=2r\,e^{\pi i s}.
\]
Thus the holomorphic projection has angular coordinate \(\pi s\):
it runs from \(1\) to \(-1\) on the boundary circle of \(\FM_2(\C)\).
Concatenating the exchange with the same construction after the labels
have been transposed gives the full circle; the single exchange is the
half-monodromy.  Its topological component changes the \(E_1\)
ordering. Since \(\mathrm{Conf}_2(\R^3)\simeq
\R^3\times(\R^3\setminus\{0\})\) has trivial fundamental group in the
ordered case and fundamental group \(\Z/2\) after quotienting by the
transposition, no braid-group \(\Z\) is assigned to the configuration
space of two points in the three-dimensional bulk.

The self-intersection $\mathfrak{S}_b = \cL \times_{\cM}^h \cL$
is fibered over $\FM_2(\C)$ with fiber the two-point OPE data.
The $(-1)$-shifted symplectic form $\omega_{-1}$ on $\mathfrak{S}_b$
restricts to the collision residue along $D_{1,2}$: it is the
pairing $\Res_{z_1 = z_2}[a(z_1) b(z_2) \, d\log(z_1 - z_2)]$
that defines the singular coefficients of the $\lambda$-bracket. Put
\(\zeta=z_1-z_2\).  Under the exchange half-monodromy with
topological transposition, \(\zeta\mapsto-\zeta\) and the two inputs
are swapped.  The logarithmic form is invariant:
\[
d\log(-\zeta)=d\log\zeta .
\]
The sign in skew-symmetry is instead the oriented endpoint sign in
Stokes' theorem on the exchange cylinder.  Projecting that Stokes
identity to the singular sector gives
\[
m_2^{\mathrm{sing}}(a,b;\zeta)
+(-1)^{\degree{a}\degree{b}}\,
\tau^\ast m_2^{\mathrm{sing}}(b,a;\zeta)
= Q(\text{exchange homotopy})
\]
up to the standard \(Q\)-exact input-correction terms displayed in
\eqref{eq:exchange-sing}.  On \(Q\)-cohomology this becomes
\[
m_2^{\mathrm{sing}}(a,b;\zeta)
=-(-1)^{\degree{a}\degree{b}}\,
\tau^\ast m_2^{\mathrm{sing}}(b,a;\zeta).
\]
The Borel transform converts \(\tau^\ast\) into
\(\lambda\mapsto-\lambda-\partial\). Hence
\(\{a_\lambda b\}=-(-1)^{\degree{a}\degree{b}}
\{b_{-\lambda-\partial}a\}\).
\end{proof}

\begin{lemma}[Chain-level exchange homotopies;
\ClaimStatusProvedHere; licensing $\alpha+\gamma$ via the oriented
exchange half-monodromy and the logarithmic chain/pro ambient
$\hypAmbientPro$]
\label{lem:PVA2_proof}
Work in the logarithmic chain/pro ambient \(\hypAmbientPro\), and fix
the oriented exchange path of Construction~\ref{const:exchange-cylinder}.
Let $H^{\mathrm{reg}}_{\mathrm{ex}}(a,b)$ and
$H^{\mathrm{sing}}_{\mathrm{ex}}(a,b;\zeta)$ be the chain homotopies obtained by
pulling the two-point weight form along the cylinder swept out by
\eqref{eq:exchange-path} and then projecting to the regular and singular sectors.
Then:
\begin{align}
\label{eq:exchange-reg}
\mu(a,b)-(-1)^{\degree{a}\degree{b}}\mu(b,a)
&=
QH^{\mathrm{reg}}_{\mathrm{ex}}(a,b)
-H^{\mathrm{reg}}_{\mathrm{ex}}(Qa,b)
-(-1)^{\degree{a}}H^{\mathrm{reg}}_{\mathrm{ex}}(a,Qb),\\[4pt]
\label{eq:exchange-sing}
m_2^{\mathrm{sing}}(a,b;\zeta)
+(-1)^{\degree{a}\degree{b}}
\,\tau^\ast m_2^{\mathrm{sing}}(b,a;\zeta)
&=
QH^{\mathrm{sing}}_{\mathrm{ex}}(a,b;\zeta)
-H^{\mathrm{sing}}_{\mathrm{ex}}(Qa,b;\zeta) \notag\\
&\quad
-(-1)^{\degree{a}}H^{\mathrm{sing}}_{\mathrm{ex}}(a,Qb;\zeta),
\end{align}
where $\tau^\ast$ is the exchange operator on singular Laurent kernels:
it exchanges the inputs, sends the collision coordinate
\(\zeta=z_1-z_2\) to \(-\zeta=z_2-z_1\), and becomes the single
vertex-algebra substitution \(\lambda\mapsto-\lambda-\partial\) after
the polynomial Borel transform and translation covariance.
\end{lemma}

\begin{proof}
Along \eqref{eq:exchange-path} one has
\[
\zeta(s)=z_1(s)-z_2(s)=2r\,e^{\pi i s}.
\]
Thus the cylinder is the oriented half-monodromy in the holomorphic
boundary circle together with the transposition of the \(E_1\)-ordered
topological coordinate. Pull the two-point logarithmic kernel to this
interval family. Because the theory is topological in the \(t\)-direction
and the logarithmic form is \((Q+d_{\mathrm{conf}})\)-closed on the
configuration space, Cartan's homotopy formula applies.  If
\(\widetilde\omega_2(a,b;s)\) is the pulled-back two-point weight form,
define
\[
\Xi_{\mathrm{ex}}(a,b):=\int_0^1 \iota_{\partial_s}\widetilde\omega_2(a,b;s)\,ds.
\]
Stokes' theorem on the oriented interval gives
\[
(Q+d)\Xi_{\mathrm{ex}}(a,b)
=
\widetilde\omega_2(a,b)\big|_{s=1}
-
\widetilde\omega_2(a,b)\big|_{s=0}.
\]
At $s=0$ we recover the original ordering. At $s=1$ the inputs are exchanged.
For the regular part $\mu$, which has degree $0$ after taking the constant term,
the exchange sign is the ordinary graded-commutative sign
$(-1)^{\degree{a}\degree{b}}$.

For the singular part the Borel-transformed $\lambda$-bracket uses
the standard vertex-algebra skew convention. The visible exchange
sign on cohomology is therefore the ordinary Koszul sign
$(-1)^{\degree{a}\degree{b}}$; the degree shift is already carried by
the bar desuspension in the source of the operation. The involution
$\tau^\ast$ records that exchanging the two points sends the collision coordinate
$\zeta=z_1-z_2$ to $-\zeta=z_2-z_1$.  Since
\[
d\log(-\zeta)=d\log\zeta,
\]
the logarithmic form contributes no extra sign.  The minus sign in the
cohomological skew identity is the oriented endpoint sign in the
Stokes formula.  After the polynomial Borel transform, translation
covariance of the second insertion converts \(\tau^\ast\) into the
standard substitution \(\lambda\mapsto-\lambda-\partial\); there is no
separate PVA operation \(\lambda\mapsto-\lambda\).
Projecting the Stokes identity to the regular and singular sectors and
using the degree-\(2\) identity \eqref{eq:repaired-n2} for input
corrections yields \eqref{eq:exchange-reg} and \eqref{eq:exchange-sing}.
\end{proof}

\begin{proposition}[Graded commutativity of the product;
\ClaimStatusProvedHere; licensing $\alpha+\gamma$ via the oriented
exchange half-monodromy and the logarithmic chain/pro ambient
$\hypAmbientPro$]
\label{prop:product-commutative}
On cohomology,
\[
[a]\cdot[b]:=[\mu(a,b)]
\]
is graded commutative:
\[
[a]\cdot[b]=(-1)^{\degree{a}\degree{b}}[b]\cdot[a].
\]
\end{proposition}

\begin{proof}
If $Qa=Qb=0$, \eqref{eq:exchange-reg} shows that
$\mu(a,b)-(-1)^{\degree{a}\degree{b}}\mu(b,a)$ is $Q$-exact. Passing to
cohomology gives the claim.
\end{proof}

\begin{proposition}[Skew-symmetry of the $\lambda$-bracket;
\ClaimStatusProvedHere; licensing $\alpha+\gamma$ via the oriented
exchange half-monodromy and the logarithmic chain/pro ambient
$\hypAmbientPro$]
\label{prop:PVA2_proof}
For cohomology classes $[a],[b]\in H^\bullet(\A,Q)$,
\begin{equation}
\label{eq:repaired-skew}
\{[a]_\lambda [b]\}
=
-(-1)^{\degree{a}\degree{b}}
\{[b]_{-\lambda-\partial}[a]\}.
\end{equation}
\end{proposition}

\begin{proof}
Take the singular part of the exchange identity \eqref{eq:exchange-sing} on
$Q$-closed representatives. The right-hand side is $Q$-exact, hence vanishes in
cohomology. Now apply the Borel transform coefficientwise. Since the Borel transform
commutes with linear combinations and with the translation operator $\partial$ on
coefficients, it converts the $\tau^\ast$-action on the Laurent kernel into the standard
polynomial substitution $\lambda\mapsto-\lambda-\partial$. This gives
\eqref{eq:repaired-skew}.
\end{proof}

\subsection{Associativity from the regular--regular sector}

\begin{proposition}[Associativity of the descended product;
\ClaimStatusProvedHere; licensing $\gamma$ via the logarithmic
chain/pro ambient $\hypAmbientPro$ and the unsuspended regular-product
convention]
\label{prop:product-associative}
The product $[a]\cdot[b]:=[\mu(a,b)]$ is associative on $H^\bullet(\A,Q)$.
\end{proposition}

\begin{proof}
Work in the unsuspended regular convention of
\cref{thm:cohomology-PVA-main}:
\[
\mu(a,b)=\mu_2^{\mathrm{reg}}(a,b;0),
\]
where \(m_2^{\mathrm{reg}}\) is the suspended bar component of the
ordinary product \(\mu_2^{\mathrm{reg}}\). Let
\(K_3^{\mathrm{reg}}(a,b,c)\) be the regular--regular constant-term
projection of the ternary homotopy \(m_3(a,b,c)\), transported by this
same desuspension convention.

Apply \(D_A^2=0\), equivalently the \(n=3\) specialization of
\eqref{eq:ainfty-relation-raw}, to \(Q\)-closed inputs \(a,b,c\).
Project to the regular--regular sector and set the collision variables
to \(0\). The input-correction terms containing \(Qa,Qb,Qc\) vanish.
In shifted \(m_2\)-coordinates the identity contains the usual
\((-1)^{\degree{a}}\) Stasheff sign. After transporting both binary
vertices through the bar desuspension to the ordinary
\(\mu_2^{\mathrm{reg}}\)-coordinates, the binary part is the ordinary
associator
\[
\mu(\mu(a,b),c)-\mu(a,\mu(b,c)).
\]
The remaining term is the \(Q\)-image of
\(K_3^{\mathrm{reg}}(a,b,c)\), up to the global sign fixed by the
chosen suspension convention. Thus
\[
\mu(\mu(a,b),c)-\mu(a,\mu(b,c))\in \operatorname{im}Q .
\]
No vanishing of \(m_3\) is assumed; \(m_3\) supplies the chain homotopy
whose regular--regular constant term kills the associator. Passing to
cohomology gives strict associativity.
\end{proof}

\subsection{The three-face Stokes theorem on $\FM_3(\C)$}
\label{subsec:three-face-stokes}

The Jacobi and Leibniz identities for the PVA bracket arise from
Stokes' theorem on the three-point configuration space $\FM_3(\C)$.
The key geometric fact is that $\partial\FM_3(\C)$ has \emph{three}
pairwise collision divisors $D_{\{1,2\}}$, $D_{\{2,3\}}$, and
$D_{\{1,3\}}$, plus a total-collision face $D_{\{1,2,3\}}$: the
first two correspond to consecutive collisions, but the third
(the non-consecutive face $D_{\{1,3\}}$) is the geometric
source of the third Jacobi term. Discarding it would break Jacobi.

The $3$-point geometry requires all three pairwise divisors:
\[
D_{\{1,2\}},\qquad D_{\{2,3\}},\qquad D_{\{1,3\}}\subset \partial\FM_3(\C).
\]
The orientation signs are
\[
\epsilon_{12}=+1,\qquad \epsilon_{23}=+1,\qquad \epsilon_{13}=-1,
\]
as computed in Appendix~\ref{app:orientations}. The minus sign on $D_{\{1,3\}}$
is exactly what is needed for the third Jacobi/Leibniz term.
In the Jacobi proof below we use the normal form
\[
J_{\lambda,\mu}(a,b,c)
:=
\{a_\lambda\{b_\mu c\}\}
-
(-1)^{(\degree{a}+1)(\degree{b}+1)}
\{b_\mu\{a_\lambda c\}\}
-
\{\{a_\lambda b\}_{\lambda+\mu}c\}.
\]
Thus the \(D_{\{1,2\}}\) contribution appears with a minus sign because
the \(\{\{a_\lambda b\}_{\lambda+\mu}c\}\) term is placed on the right
side of \eqref{eq:repaired-jacobi}; the \(D_{\{1,3\}}\) contribution
appears with a minus sign from the incidence orientation
\(\epsilon_{13}=-1\) together with the shifted Koszul transposition.

\begin{lemma}[Three-face singular factorization;
\ClaimStatusProvedHere; licensing $\gamma$ via the logarithmic
chain/pro ambient $\hypAmbientPro$, the oriented $\FM_3$ incidence
convention, and the divided-power Borel transform]
\label{lem:PVA3_proof}
Let $\Omega_3^{\mathrm{sing}\text{-}\mathrm{sing}}(a,b,c)$ be the doubly singular sector of the
three-point logarithmic form. In the Jacobiator normalization above,
the oriented pairwise boundary contributions are:
\begin{align}
\label{eq:res23-jac}
\mathfrak J_{23}
&=
\beta(a,\beta(b,c))\ \leadsto\ \{a_\lambda\{b_\mu c\}\},\\[4pt]
\label{eq:res13-jac}
\mathfrak J_{13}
&=
-\,(-1)^{(\degree{a}+1)(\degree{b}+1)}
\,\beta(b,\beta(a,c))
\ \leadsto\
-\,(-1)^{(\degree{a}+1)(\degree{b}+1)}
\{b_\mu\{a_\lambda c\}\},\\[4pt]
\label{eq:res12-jac}
\mathfrak J_{12}
&=
-\,\beta(\beta(a,b),c)\ \leadsto\ -\,\{\{a_\lambda b\}_{\lambda+\mu}c\}.
\end{align}
The symbol ``$\leadsto$'' means: apply the Borel transform in the relevant singular
variables and substitute the outer fused parameter by the sum of the inner parameters.
\end{lemma}

\begin{proof}
Near each divisor one chooses standard FM blow-up coordinates.

\smallskip
\noindent\textbf{Face $D_{\{1,2\}}$.}
Write
\[
\zeta_{12}=z_1-z_2=\varepsilon_{12}e^{i\theta_{12}},
\qquad
u_{12}=\frac{z_1+z_2}{2}-z_3.
\]
The product decomposition in Definition~\ref{def:log-SC-algebra} implies that, up to terms regular in $\varepsilon_{12}$,
the three-point form splits as
\[
d\log\varepsilon_{12}\wedge
\omega_2(a,b;\zeta_{12})
\wedge
\omega_2\bigl(m_2(a,b),c;u_{12}\bigr).
\]
Taking the singular residue in the inner and outer channels yields the iterated singular
composition $\beta(\beta(a,b),c)$. The outer fused collision parameter
is the sum of the two incoming polynomial parameters, hence
\(\lambda+\mu\). In the Jacobiator normalization this is the term
\(-\{\{a_\lambda b\}_{\lambda+\mu}c\}\), because
\eqref{eq:repaired-jacobi} places it on the right-hand side.

\smallskip
\noindent\textbf{Face $D_{\{2,3\}}$.}
Similarly, with
\[
\zeta_{23}=z_2-z_3=\varepsilon_{23}e^{i\theta_{23}},
\qquad
u_{23}=z_1-\frac{z_2+z_3}{2},
\]
factorization gives the ordered composition $\beta(a,\beta(b,c))$, which becomes
$\{a_\lambda\{b_\mu c\}\}$ after Borel transform.

\smallskip
\noindent\textbf{Face $D_{\{1,3\}}$.}
Here
\[
\zeta_{13}=z_1-z_3=\varepsilon_{13}e^{i\theta_{13}},
\qquad
u_{13}=z_2-\frac{z_1+z_3}{2}.
\]
The residue again factors into an inner and an outer binary piece, but now the
cluster $\{1,3\}$ has to be moved past the external input $b$. Because the singular
binary piece has degree $-1$, the transposition contributes the shifted Koszul sign
$(-1)^{(\degree{a}+1)(\degree{b}+1)}$. Combining this with the orientation sign
$\epsilon_{13}=-1$ produces the oriented contribution
\[
-\,(-1)^{(\degree{a}+1)(\degree{b}+1)}
\{b_\mu\{a_\lambda c\}\}.
\]
This does not contradict the non-consecutive vanishing statement for
planar tree-level weight forms: the present form
\(\Omega_3^{\mathrm{sing}\text{-}\mathrm{sing}}\) is the dressed
doubly singular PVA form, and
Remark~\ref{rem:non-consecutive-scope} says that its
\(D_{\{1,3\}}\)-residue is generically non-zero.
\end{proof}

\begin{theorem}[Jacobi identity;
\ClaimStatusProvedHere; licensing $\gamma$ via the logarithmic
chain/pro ambient $\hypAmbientPro$, the oriented $\FM_3$ incidence
convention, and Arnold--Orlik--Solomon corner cancellation]
\label{thm:Jacobi}
For all cohomology classes $[a],[b],[c]\in H^\bullet(\A,Q)$,
\begin{equation}
\label{eq:repaired-jacobi}
\{[a]_\lambda\{[b]_\mu[c]\}\}
-
(-1)^{(\degree{a}+1)(\degree{b}+1)}
\{[b]_\mu\{[a]_\lambda[c]\}\}
=
\{\{[a]_\lambda[b]\}_{\lambda+\mu}[c]\}.
\end{equation}
\end{theorem}

\begin{proof}
Apply Stokes' theorem to the doubly singular three-point form.  Let
\(K_3^{\mathrm{Jac}}(a,b,c)\) be the doubly singular component of the
interior \(m_3\)-homotopy.  After inserting the incidence orientations
and projecting to the singular--singular sector, Stokes gives
\[
J_{\lambda,\mu}(a,b,c)\in\operatorname{im}Q,
\qquad
J_{\lambda,\mu}:=\mathfrak J_{23}+\mathfrak J_{13}+\mathfrak J_{12}.
\]
Equivalently, the difference between \(J_{\lambda,\mu}(a,b,c)\) and
zero is \(QK_3^{\mathrm{Jac}}(a,b,c)\), up to the global sign fixed by
the suspension convention. The total-collision divisor contributes
only this \(Q\)-boundary, so it disappears in cohomology.

Lemma~\ref{lem:PVA3_proof} identifies the three oriented boundary
contributions with the three terms of the Jacobiator:
\[
J_{\lambda,\mu}(a,b,c)
=
\{a_\lambda\{b_\mu c\}\}
-
(-1)^{(\degree{a}+1)(\degree{b}+1)}
\{b_\mu\{a_\lambda c\}\}
-
\{\{a_\lambda b\}_{\lambda+\mu}c\}.
\]
The form identity at the corner is the Arnold relation
\[
\omega_{12}\wedge\omega_{23}
\;+\;
\omega_{23}\wedge\omega_{31}
\;+\;
\omega_{31}\wedge\omega_{12}=0,
\qquad
\omega_{ij}=d\log(z_i-z_j).
\]
The three summands are the logarithmic forms underlying
\(\mathfrak J_{12}\), \(\mathfrak J_{23}\), and
\(\mathfrak J_{13}\), with the incidence and Koszul signs already
included in the \(\mathfrak J\)-notation. Thus the Arnold relation is
the printed Jacobiator \(J_{\lambda,\mu}\).
Arnold--Orlik--Solomon cancellation at codimension-$2$ corners
(Theorem~\ref{thm:Arnold_full_proof}, Appendix~\ref{app:FM_Stokes})
ensures that no extra corner contribution survives.
Hence \(J_{\lambda,\mu}\) is \(Q\)-exact; on cohomology it vanishes.
\end{proof}

\begin{remark}[Conditional Steinberg reading of the Jacobi identity]
\label{rem:jacobi-lagrangian}
The three boundary faces $D_{\{1,2\}},D_{\{2,3\}},D_{\{1,3\}}\subset
\partial\FM_3(\C)$ appearing in the Stokes argument admit a Steinberg
interpretation only after the comparison data of
Proposition~\ref{prop:PVA-from-symplectic} have been supplied.  In the
bar/Koszul model of the self-intersection \(\mathfrak S_b\), they are
the three pairwise projections of a triple Lagrangian correspondence.
When the comparison map \(\chi\) is an isomorphism on the
Koszul-effective locus \(\effKoszul\), the FM Jacobi identity on
boundary cohomology is the image of the associativity of this
correspondence convolution.  Without that comparison, the convolution
statement remains a statement on the Steinberg bar/Koszul shadow; the
boundary PVA Jacobi identity is still the direct FM theorem
\ref{thm:Jacobi}.
\end{remark}

\begin{proposition}[Conditional Steinberg convolution interpretation of
Jacobi; \ClaimStatusConditional; licensing
$\alpha+\beta+\gamma+\varepsilon$ via boundary chart $b$, the
comparison map $\chi$, the derived shifted-symplectic ambient, and
$\effKoszul$]
\label{prop:jacobi-lagrangian-convolution}
Assume the hypotheses of Proposition~\textup{\ref{prop:PVA-from-symplectic}}.
Let
\[
\chi\colon H^\bullet(\barB(\A))\longrightarrow H^\bullet(\A,Q)
\]
be the bar-to-boundary comparison map, and assume that \(\chi\) is an
isomorphism on the Koszul-effective comparison locus \(\effKoszul\).
Then the signed Jacobiator \(J_{\lambda,\mu}\) of
Theorem~\textup{\ref{thm:Jacobi}} is the \(\chi\)-image of the
associator of Lagrangian triple convolution in \(\Steinb_b\). Hence,
on this comparison locus, the boundary PVA Jacobi identity agrees with
associativity of Steinberg convolution.  Without the isomorphism
hypothesis on \(\chi\), this is only a Steinberg bar/Koszul-shadow
interpretation and not a proof of the boundary PVA identity.
\end{proposition}

\begin{proof}
The direct boundary identity has already been proved in
Theorem~\ref{thm:Jacobi} by the signed three-face Stokes argument.  It
remains only to identify its Steinberg shadow under the comparison
datum.

In the bar/Koszul model of the derived self-intersection, the three
pairwise collision faces
\[
D_{\{1,2\}},\qquad D_{\{2,3\}},\qquad D_{\{1,3\}}
\]
correspond to the three pair-projections of the triple fibre product
\[
\mathscr L_b\times_{\mathscr M}\mathscr L_b\times_{\mathscr M}\mathscr L_b .
\]
Indeed, under the identification of \(\barB(\A)\) with the formal
bar/Koszul model of the derived self-intersection
\textup{(}Theorem~\textup{\ref{thm:bar-is-self-intersection}}\textup{)},
the self-intersection
\[
\Steinb_b=\mathscr L_b\times_{\mathscr M}\mathscr L_b
\]
acts as a correspondence from \(\mathscr L_b\) to itself, and the
triple fibre product admits pair-projections
\[
p_{12},\quad p_{23},\quad p_{13}
\;\colon\;
\Steinb_b \times_{\mathscr{L}_b} \Steinb_b
\;\longrightarrow\;
\Steinb_b,
\]
one for each pairwise composition.  With the signed Jacobiator
normalization of Theorem~\ref{thm:Jacobi}, these projections correspond
to
\begin{itemize}
\item $D_{\{2,3\}}$: compose via $p_{23}$, yielding
 $\{a_\lambda \{b_\mu c\}\}$;
\item $D_{\{1,3\}}$: compose via $p_{13}$, yielding
 $-(-1)^{(\degree{a}+1)(\degree{b}+1)}
 \{b_\mu \{a_\lambda c\}\}$;
\item $D_{\{1,2\}}$: compose via $p_{12}$, yielding
 $-\{\{a_\lambda b\}_{\lambda+\mu} c\}$.
\end{itemize}
Thus the Steinberg convolution associator maps to the signed sum
\[
J_{\lambda,\mu}
=
\mathfrak J_{23}+\mathfrak J_{13}+\mathfrak J_{12}.
\]
Associativity of derived correspondence composition is the canonical
equivalence between the two parenthesisations of
\[
(\mathscr L_b\times_{\mathscr M}\mathscr L_b)
\times_{\mathscr L_b}
(\mathscr L_b\times_{\mathscr M}\mathscr L_b),
\]
and the Arnold--Orlik--Solomon corner cancellation is its logarithmic
de Rham representative.  Applying the comparison map \(\chi\) sends
this Steinberg associator to the boundary Jacobiator.  If \(\chi\) is
an isomorphism on \(\effKoszul\), the Steinberg associativity statement
is equivalent to the boundary Jacobi identity.  If not, the conclusion
remains only on the bar/Koszul shadow, exactly as in
Proposition~\ref{prop:PVA-from-symplectic}.
\end{proof}

\subsection{Leibniz from the mixed regular--singular sector}

The Jacobi identity was derived from the doubly singular sector of
the three-point form (both binary channels projected to singular
parts). The Leibniz rule comes from the \emph{mixed} sector: one
binary channel projected to the singular part (producing a
$\lambda$-bracket) and the other to the regular part at collision
parameter~$0$ (producing a product). The same three-face Stokes
argument applies, with the same oriented \(\FM_3\) incidence
convention.  We use the signed Leibniz defect
\[
\mathcal L_\lambda(a;b,c)
:=
\{a_\lambda(b\cdot c)\}
-
\{a_\lambda b\}\cdot c
-
(-1)^{(\degree{a}+1)\degree{b}}\,
b\cdot\{a_\lambda c\}.
\]

\begin{lemma}[Mixed three-face factorization;
\ClaimStatusProvedHere; licensing $\gamma$ via the logarithmic
chain/pro ambient $\hypAmbientPro$, the oriented $\FM_3$ incidence
convention, the unsuspended regular product, and the divided-power
Borel transform]
\label{lem:three-face-leibniz}
Let $\Omega_3^{\mathrm{mix}}(a,b,c)$ be the sector in which one binary channel is
projected to the singular part and the other to the regular part, with the outer regular
part evaluated at collision parameter $0$. In the Leibniz-defect
normalization above, the oriented pairwise contributions are:
\begin{align}
\label{eq:res23-leibniz}
\mathfrak L_{23} &\leadsto \{a_\lambda (b\cdot c)\},\\
\label{eq:res12-leibniz}
\mathfrak L_{12} &\leadsto -\,\{a_\lambda b\}\cdot c,\\
\label{eq:res13-leibniz}
\mathfrak L_{13} &\leadsto
-\,(-1)^{(\degree{a}+1)\degree{b}}\,
b\cdot\{a_\lambda c\}.
\end{align}
\end{lemma}

\begin{proof}
The local-coordinate analysis is the same as in Lemma~\ref{lem:PVA3_proof}, but now
one takes one singular residue and one regular constant term.

Near $D_{\{2,3\}}$, the inner collision of $b$ and $c$ is projected to the regular
constant term, producing the product $bc$, while the outer channel is singular, giving
$\{a_\lambda(bc)\}$.

Near $D_{\{1,2\}}$, the inner collision is singular and gives $\{a_\lambda b\}$; the
outer channel is regular at parameter $0$, giving the product with
$c$. This is the first right-hand term of the Leibniz rule and hence
appears with the minus sign in \(\mathcal L_\lambda\).

Near $D_{\{1,3\}}$, the inner singular collision gives $\{a_\lambda c\}$ and the
outer regular channel multiplies the result by $b$. To move the degree-$(-1)$ bracket
past the element $b$ one picks up the graded Poisson sign
$(-1)^{(\degree{a}+1)\degree{b}}$; in the signed defect this is the
second right-hand term and therefore appears as
\[
-\,(-1)^{(\degree{a}+1)\degree{b}}\,b\cdot\{a_\lambda c\}.
\]
The mixed sector contains exactly one singular channel, so the
singular--singular Wick integral term of a quantum vertex algebra does
not appear in this classical PVA descent.
\end{proof}

\begin{proposition}[Leibniz rule in cohomology;
\ClaimStatusProvedHere; licensing $\gamma$ via the logarithmic
chain/pro ambient $\hypAmbientPro$, the oriented $\FM_3$ incidence
convention, the unsuspended regular product, and the divided-power
Borel transform]
\label{prop:PVA3_proof}
\label{prop:Leibniz}
For all cohomology classes $[a],[b],[c]\in H^\bullet(\A,Q)$,
\begin{equation}
\label{eq:repaired-leibniz}
\{[a]_\lambda([b]\cdot[c])\}
=
\{[a]_\lambda[b]\}\cdot[c]
+
(-1)^{(\degree{a}+1)\degree{b}}\,[b]\cdot\{[a]_\lambda[c]\}.
\end{equation}
\end{proposition}

\begin{proof}
Apply Stokes' theorem to $\Omega_3^{\mathrm{mix}}(a,b,c)$. Let
\(K_3^{\mathrm{Leib}}(a;b,c)\) be the mixed regular--singular
component of the interior \(m_3\)-homotopy. After inserting the
incidence orientations and projecting to the mixed sector, Stokes gives
\[
\mathcal L_\lambda(a;b,c)\in \operatorname{im}Q.
\]
Equivalently,
\(\mathcal L_\lambda(a;b,c)=QK_3^{\mathrm{Leib}}(a;b,c)\), up to the
global sign fixed by the suspension convention. The total-collision
face contributes only this \(Q\)-boundary, so it disappears in
cohomology. Lemma~\ref{lem:three-face-leibniz} identifies the three
pairwise oriented contributions with the three summands of
\(\mathcal L_\lambda\), and the same Arnold--Orlik--Solomon corner
cancellation used in the Jacobi proof ensures that no additional
corner term remains. Hence the Leibniz defect is \(Q\)-exact; on
cohomology it vanishes, which is \eqref{eq:repaired-leibniz}.
\end{proof}

\subsection{Vacuum and completion of the PVA proof}

With sesquilinearity, skew-symmetry, associativity, Jacobi, and
Leibniz now established, it remains to verify the vacuum axiom
and assemble the complete PVA structure.

\begin{proposition}[Vacuum/unit axiom;
\ClaimStatusProvedHere; licensing $\gamma$ via the strict unital
$A_\infty$-chiral structure in the logarithmic chain/pro ambient
$\hypAmbientPro$, the unsuspended regular product, the
regular/singular projection, and the divided-power Borel transform]
\label{prop:PVA4_proof}
Assume the strict unit \(\mathbf 1\in\A^0\) satisfies
\(Q\mathbf 1=0\), \(\partial\mathbf 1=0\), the binary unit relation
\eqref{eq:unit-m2}, and the higher-unit relation
\eqref{eq:unit-higher}.  Then the class
$[\mathbf 1]\in H^\bullet(\A,Q)$ is a vacuum vector for the descended
PVA:
\[
[\mathbf 1]\cdot[a]=[a]\cdot[\mathbf 1]=[a],
\qquad
\{[\mathbf 1]_\lambda[a]\}=0,
\qquad
\{[a]_\lambda[\mathbf 1]\}=0,
\qquad
\partial[\mathbf 1]=0.
\]
\end{proposition}

\begin{proof}
The suspended bar operation satisfies \eqref{eq:unit-m2}.  After the
bar desuspension used throughout the product proof, the descended
product is represented by the ordinary unsuspended regular product
\(\mu_2^{\mathrm{reg}}(-,-;0)\), for which the strict unit identities
are
\[
\mu_2^{\mathrm{reg}}(\mathbf 1,a;0)=a,\qquad
\mu_2^{\mathrm{reg}}(a,\mathbf 1;0)=a .
\]
The corresponding two binary kernels are constant in the collision
coordinate.  Hence their singular projections vanish:
\[
m_2^{\mathrm{sing}}(\mathbf 1,a)=0,
\qquad
m_2^{\mathrm{sing}}(a,\mathbf 1)=0.
\]
After the divided-power Borel transform this gives
\(\{[\mathbf 1]_\lambda[a]\}=0\) and
\(\{[a]_\lambda[\mathbf 1]\}=0\).  The higher-unit axiom
\eqref{eq:unit-higher} says that no \(m_k\), \(k\ne2\), contributes an
extra boundary term with a unit insertion.  Finally
\(Q\mathbf 1=0\) makes \([\mathbf 1]\) a cohomology class and
\(\partial\mathbf 1=0\) gives \(\partial[\mathbf 1]=0\).
\end{proof}

\begin{corollary}[Shifted PVA structure on cohomology;
\ClaimStatusProvedHere; licensing $\alpha+\gamma$ via the oriented
exchange half-monodromy, the logarithmic chain/pro ambient
$\hypAmbientPro$, the oriented $\FM_3$ incidence convention, the
strict unital $A_\infty$-chiral structure, the unsuspended regular
product, and the divided-power Borel transform]
\label{cor:PVA-structure}
The cohomology $H^\bullet(\A,Q)$ equipped with
\[
[a]\cdot[b]:=[\mu(a,b)],
\qquad
\{[a]_\lambda[b]\}:=\Bigl[\sum_{n\ge0}\frac{\lambda^n}{n!}\,a_{(n)}b\Bigr],
\qquad
\partial,
\qquad
[\mathbf 1]
\]
is a unital, associative, graded-commutative differential algebra together with a
degree-$(-1)$ $\lambda$-bracket satisfying sesquilinearity, ordinary
vertex-algebra skew-symmetry, Jacobi, and Leibniz with the shifted
Koszul signs. Equivalently, it is a $(-1)$-shifted Poisson vertex algebra.
\end{corollary}

\begin{proof}
Combine Proposition~\ref{prop:PVA1_proof},
Proposition~\ref{prop:product-commutative},
Proposition~\ref{prop:product-associative},
Proposition~\ref{prop:PVA2_proof},
Theorem~\ref{thm:Jacobi},
Proposition~\ref{prop:PVA3_proof},
and Proposition~\ref{prop:PVA4_proof}.
\end{proof}

\subsection{Conditional higher vanishing}

The previous subsections close the PVA layer itself. They do not
prove that the whole $A_\infty$ structure is formal. Contractibility
of the ordered topological configuration space is not by itself a
chain-level nullhomotopy of the operations: the compactified
configuration space has collision faces, and those faces are precisely
where the Stasheff terms live. Higher operations vanish on cohomology
only after a compatible topological bounding-chain datum is supplied.

\begin{lemma}[Open ordered topological contraction;
\ClaimStatusProvedHere; licensing $\gamma$ via the ordered
topological chain ambient inside $\hypAmbientPro$ and the translation
quotient]
\label{lem:topological-contraction}
Fix $k\ge3$. The ordered configuration space modulo translation,
\[
\Conf_k^{<}(\R)/\mathrm{transl}
\;=\;
\{(s_1,\ldots,s_{k-1})\in\R^{k-1}_{>0}\},
\qquad
s_j:=t_{j+1}-t_j,
\]
is diffeomorphic to the open positive orthant $\R^{k-1}_{>0}$ and therefore contractible.
In particular $H_j(\Conf_k^{<}(\R))=0$ for all $j\ge1$. The linear retraction
\begin{equation}
\label{eq:affine-retraction}
\rho_u\colon (s_1,\ldots,s_{k-1})\;\longmapsto\;
(1-u)(s_1,\ldots,s_{k-1})+u(1,\ldots,1),
\qquad u\in[0,1],
\end{equation}
contracts $\Conf_k^{<}(\R)/\mathrm{transl}$ to the interior basepoint
$(1,\ldots,1)$.
This is a statement about the open ordered topological factor only; it
does not by itself construct a relative bounding chain in the
compactified operadic pair.
\end{lemma}

\begin{proof}
The map $(t_1,\ldots,t_k)\mapsto (t_2-t_1,\ldots,t_k-t_{k-1})$ identifies the ordered
configuration modulo translation with $\R^{k-1}_{>0}$. This is star-shaped about any interior
point, hence contractible; in particular, the straight-line homotopy
$\rho_u$ in~\eqref{eq:affine-retraction} is smooth and stays in
$\R^{k-1}_{>0}$ for all $u\in[0,1]$.
\end{proof}

\begin{definition}[Compatible topological nullhomotopies]
\label{def:compatible-topological-nullhomotopy}
A logarithmic $\SCchtop$-algebra has \emph{compatible topological
nullhomotopies} if, for every $k\ge3$, the topological chain
$\omega_k^{\mathrm{top}}$ defining the $k$-ary operation admits a
chain $\Gamma_{k-1}$ in the compactified ordered topological factor,
relative to its operadic collision boundary, such that
\[
\partial_{\mathrm{rel}}\Gamma_{k-1}=\omega_k^{\mathrm{top}},
\]
and the product decomposition
$\widetilde\omega_k=\omega_k^{\mathrm{hol}}\wedge\omega_k^{\mathrm{top}}$
extends over $\Gamma_{k-1}$ without creating new holomorphic
boundary residues. Equivalently, the topological boundary of
$\Gamma_{k-1}$ is identified, through the
$C_\ast(W(\SCchtop))$-algebra differential, with the output
$Q$-boundary of the corresponding $k$-ary operation.
\end{definition}

\begin{proposition}[Higher operations are \(Q\)-exact under compatible topological
nullhomotopies; \ClaimStatusProvedHere; licensing $\gamma+\varepsilon$
via the logarithmic chain/pro ambient $\hypAmbientPro$, the
relative compactified ordered topological factor, compatible
topological nullhomotopies in
Definition~\textup{\ref{def:compatible-topological-nullhomotopy}}, and
the factorisation
\(\widetilde\omega_k=\omega_k^{\mathrm{hol}}\wedge\omega_k^{\mathrm{top}}\)]
\label{prop:m3_vanish}
Let $(\cA_{\mathrm{ch}}, \cA_{\mathrm{top}})$ be a logarithmic $\SCchtop$-algebra (Definition~\ref{def:log-SC-algebra}).
Assume it has compatible topological nullhomotopies in the sense of
Definition~\ref{def:compatible-topological-nullhomotopy}.
For every $k\ge3$ and all $Q$-closed elements $a_1,\ldots,a_k\in\A$,
there exists $h_{k-1}(a_1,\ldots,a_k)\in\A((\zeta_1))\cdots((\zeta_{k-1}))$ such that
\begin{equation}
\label{eq:higher-Q-exact}
m_k(a_1,\ldots,a_k)
\;=\;
Q\bigl(h_{k-1}(a_1,\ldots,a_k)\bigr).
\end{equation}
In particular, under this compatible-nullhomotopy hypothesis, the
direct operation \(m_k\) represents zero on \(H^\bullet(\A,Q)\) for
every \(k\ge3\).  Without the supplied relative bounding chains, no
higher-operation vanishing statement is asserted.
\end{proposition}

\begin{proof}
The proof has four steps.

\smallskip
\noindent\textbf{Step~1 (Factorisation of the weight form).}
By Definition~\ref{def:log-SC-algebra}, the weight form defining $m_k$ factors as
\begin{equation}
\label{eq:weight-factor}
\widetilde\omega_k(a_1,\ldots,a_k)
\;=\;
\omega_k^{\mathrm{hol}}(a_1,\ldots,a_k)
\;\wedge\;
\omega_k^{\mathrm{top}},
\end{equation}
where $\omega_k^{\mathrm{hol}}\in\Omega^\bullet_{\log}(\FM_k(\C))\otimes\A^{\otimes k}$
is the holomorphic logarithmic kernel and
$\omega_k^{\mathrm{top}}\in C_{k-2}(\Conf_k^{<}(\R)/\mathrm{transl})$ is the topological
fundamental chain (the space $\Conf_k^{<}(\R)/\mathrm{transl}\cong\R^{k-1}_{>0}$
has dimension $k-1$; the degree is $k-2$ rather than $k-1$ because the operadic
desuspension $s^{-1}$ in the bar construction shifts the topological degree by~$-1$).
The operation is
\[
m_k(a_1,\ldots,a_k)
\;=\;
\int_{\FM_k(\C)\,\times\,\Conf_k^{<}(\R)}
\widetilde\omega_k(a_1,\ldots,a_k).
\]
For $k\ge3$, the topological chain $\omega_k^{\mathrm{top}}$ has degree $k-2\ge1$.

\smallskip
\noindent\textbf{Step~2 (Topological bounding chain).}
The compatible-nullhomotopy hypothesis supplies a relative
$(k{-}1)$-chain $\Gamma_{k-1}$ satisfying
\begin{equation}
\label{eq:bounding-chain}
\partial_{\mathrm{rel}}\Gamma_{k-1}=\omega_k^{\mathrm{top}}.
\end{equation}
Ordinary contractibility of the open orthant
(Lemma~\ref{lem:topological-contraction}) is only the homological
background for this hypothesis; it does not supply the required
relative chain with operadic boundary compatibility.

\smallskip
\noindent\textbf{Step~3 (Constructing the homotopy).}
Define
\begin{equation}
\label{eq:homotopy-def}
h_{k-1}(a_1,\ldots,a_k)
\;:=\;
\int_{\FM_k(\C)\,\times\,\Gamma_{k-1}}
\omega_k^{\mathrm{hol}}(a_1,\ldots,a_k)
\;\wedge\;
\mathbf{1}_{\Gamma_{k-1}}.
\end{equation}
The factorisation~\eqref{eq:weight-factor} shows that the holomorphic kernel
$\omega_k^{\mathrm{hol}}$ is \emph{untouched} by the topological contraction.
Thus the logarithmic-form, residue, and corner-cancellation properties of the
integrand remain exactly those supplied by Theorem~\ref{thm:FM-calculus}; only
the topological nullhomotopy chain is changed.

\smallskip
\noindent\textbf{Step~4 (Topological boundary and the $Q$-exactness identity).}
Use the chain-differential compatibility built into the
$C_\ast(W(\SCchtop))$-algebra structure. For a factored operation, the
internal differential $Q$ on the output is intertwined with the chain
differential on the topological input. Since the inputs are $Q$-closed,
only the output differential survives, and Stokes in the topological
factor gives
\begin{align*}
Q\bigl(h_{k-1}(a_1,\ldots,a_k)\bigr)
&=
\int_{\FM_k(\C)\,\times\,\partial_{\mathrm{rel}}\Gamma_{k-1}}
\omega_k^{\mathrm{hol}}(a_1,\ldots,a_k)
\;\wedge\;\mathbf{1}_{\partial_{\mathrm{rel}}\Gamma_{k-1}}
\\[4pt]
&=
\int_{\FM_k(\C)\,\times\,\omega_k^{\mathrm{top}}}
\widetilde\omega_k(a_1,\ldots,a_k)
\\[4pt]
&=
m_k(a_1,\ldots,a_k).
\end{align*}
The first equality uses the $C_\ast(W(\SCchtop))$ chain-map relation for the
topological factor together with~\eqref{eq:bounding-chain}; the second simply
re-assembles the factored integrand.

This establishes~\eqref{eq:higher-Q-exact}. Since $m_k(a_1,\ldots,a_k)$ is
$Q$-exact whenever all inputs are $Q$-closed, it represents the zero class in
$H^\bullet(\A,Q)$.
\end{proof}

\begin{remark}[Why $k=2$ escapes]
\label{rem:why-k2-escapes}
For $k=2$, the topological chain $\omega_2^{\mathrm{top}}\in
C_0(\Conf_2^{<}(\R)/\mathrm{transl})$ is a $0$-chain, i.e.\ a point.
Since $H_0\neq 0$, this $0$-chain is \emph{not} a boundary, so the
nullhomotopy criterion cannot kill it. This is why $m_2$ persists on
cohomology and gives rise to the product and $\lambda$-bracket.
\end{remark}

\begin{remark}[Chain-level vanishing versus transferred operations]
\label{rem:chain-vs-transferred}
Under the compatible-nullhomotopy hypothesis,
Proposition~\ref{prop:m3_vanish} shows that $m_k$ is
$Q$-exact on $Q$-closed inputs: $m_k|_{H^{\otimes k}}$
represents zero in $H^\bullet(\cA, Q)$. This is \emph{not}
the same as the vanishing of the transferred $\Ainf$
operations $m_k^H$ on cohomology from the Kadeishvili--Merkulov
transfer theorem
(Theorem~\ref{thm:resolvent-principle}). The transferred
$m_k^H\colon H^{\otimes k} \to H$ is defined via the tree
formula involving the contracting homotopy~$h$, and
measures the \emph{excess Tor} in the derived
self-intersection $\cL_\cA \times_{\Mvac} \cL_\cA$
(Corollary~\ref{cor:koszul-resolvent}).

For Koszul algebras equipped with these nullhomotopies, direct
$Q$-exactness and algebraic formality $m_k^H = 0$
(Corollary~\ref{cor:koszul-resolvent}) are consistent.
For algebras with non-formal Swiss-cheese structure (e.g.\ Virasoro at
$c \neq 0$, after DS reduction), the transferred $m_k^H$ may be
non-zero even when a direct operation has been killed on chosen
representatives: $m_k^H$ involves compositions through
non-cohomological intermediate states (via~$h$), not only direct
restriction to~$H$. Without compatible nullhomotopies, no
chain-level vanishing statement is asserted.
\end{remark}

\begin{example}[Explicit contraction at $k=3$]
\label{ex:k3-contraction}
At degree $k=3$, modding out translation via $s_1=t_2-t_1>0$, $s_2=t_3-t_2>0$
gives
\[
\Conf_3^{<}(\R)/\mathrm{transl}
\;\cong\;
\R^2_{>0}
\;=\;
\{(s_1,s_2)\,:\,s_1>0,\;s_2>0\}.
\]
Suppose the topological fundamental $1$-chain (the ``time-ordering
contour'') is represented by a compact relative cycle
$\gamma\subset\R^2_{>0}$ and is equipped with the compatible
nullhomotopy required by
Definition~\ref{def:compatible-topological-nullhomotopy}. The
interior retraction~\eqref{eq:affine-retraction} is then a model for
the background contraction:
\[
(s_1,s_2,u)\;\longmapsto\;
\bigl((1-u)\,s_1+u,\;(1-u)\,s_2+u\bigr),
\qquad u\in[0,1].
\]
The supplied nullhomotopy gives a $2$-chain
$\Gamma_2\subset\R^2_{>0}$ with
$\partial_{\mathrm{rel}}\Gamma_2=\gamma$.
The homotopy is therefore
\[
h_2(a,b,c)
\;=\;
\int_{\FM_3(\C)}
\omega_3^{\mathrm{hol}}(a,b,c)
\;\cdot\;
\int_{\Gamma_2}ds_1\,ds_2,
\]
which is finite because the FM-regularised holomorphic kernel has at worst logarithmic
singularities by Theorem~\ref{thm:FM-calculus} and the topological cone $\Gamma_2$ has finite volume
for any compactly supported representative of $\gamma$.
\end{example}

\begin{remark}[Role of the factorisation hypothesis]
\label{rem:factorisation-role}
The factorisation
$\widetilde\omega_k = \omega_k^{\mathrm{hol}}\wedge\omega_k^{\mathrm{top}}$
from Definition~\ref{def:log-SC-algebra} is the compatibility input
that makes the nullhomotopy criterion meaningful. Without
this factorisation, the topological nullhomotopy
could in principle interact with the holomorphic singular structure, creating new
poles or invalidating the FM regularisation. The factorisation ensures that
the bounding chain $\Gamma_{k-1}$ acts only in the topological fibre and leaves the
holomorphic kernel, and all its boundary residues, invariant. This is the
precise sense in which an independently supplied ordered-topological
nullhomotopy implies that the higher operations are cohomologically
trivial.
\end{remark}

\begin{remark}[The role of $d'=1$]
\label{rem:d-prime-role}
Proposition~\ref{prop:m3_vanish} is a $d'=1$ criterion (one
topological direction) with an added compatible-nullhomotopy
hypothesis.
For $d'\ge2$, the Gaiotto--Kulp--Wu formality theorem~\cite{GKW25} implies that
higher operations vanish already at the chain level (after one-loop renormalisation).
For $d'=1$, they persist at the chain level (the Landau--Ginzburg model has a
genuine chain-level $m_3$ (Example~\ref{ex:m3_LG_cubic})); they
become $Q$-exact only on loci where the compatible topological
nullhomotopies exist. The PVA descent itself uses only the binary
operation and the $n=3$ Stokes identities, not unconditional
higher-operation formality.
\end{remark}

\begin{remark}[PVA descent as further projection from $\Sym^c$ to classical]
\label{rem:pva-classical-projection}%
\index{PVA descent!projection hierarchy}%
\index{bar complex!projection to PVA}%
The PVA on $H^\bullet(\cA, Q)$ sits at the end of a three-stage
projection chain:
\[
\barB^{\mathrm{ord}}(\cA)
\;\xrightarrow{\;\mathrm{av}\;}
\barBch(\cA)
\;\xrightarrow{\;H^\bullet\;}
H^\bullet(\cA,Q)
\;\xrightarrow{\;\hbar \to 0\;}
\text{PVA}.
\]
The first step is $\Sigma_n$-coinvariant averaging (discarding the
$E_1$ ordering, losing the $R$-matrix). The second is passage to
cohomology and then to the binary classical shadow: higher
$A_\infty$ operations $m_{k \ge 3}$ are not part of the PVA data,
and vanish only when Proposition~\ref{prop:m3_vanish}'s
nullhomotopy hypothesis is present. The third
is the classical limit: the surviving binary operation $m_2$ splits
into a commutative product $\mu$ (from the regular part) and a
$\lambda$-bracket (from the singular part), assembling into a
$(-1)$-shifted Poisson vertex algebra. In the language of Volume~I,
this chain strips successively the $E_1$ data ($R$-matrix, braiding),
the homotopy-theoretic data (algebraic depth $d_{\mathrm{alg}}$),
and the non-commutativity, retaining only the arithmetic depth
$d_{\mathrm{arith}}$ (the $\lambda$-bracket's residue content).
\end{remark}

\begin{remark}[PVA descent in the three-stage projection hierarchy]
\label{rem:pva-three-stage-vol1}%
\index{PVA descent!three-stage projection}%
\index{bar complex!projection to classical PVA}%
The projection chain of
Remark~\ref{rem:pva-classical-projection} terminates at
$\Sym^c$, the symmetric-chiral level. The PVA is one step
further: it is the passage from $\Sym^c(\cA)$ (the
quantum vertex algebra on cohomology, retaining the
non-commutative normal-ordered product) to the classical
commutative algebra with $\lambda$-bracket. In Volume~I's
shadow metric formulation, the PVA shadow corresponds to
$Q_L^{\mathrm{cl}}(t) = (2\kappa + 3\alpha t)^2$: the
critical discriminant $\Delta = 8\kappa S_4$ vanishes on the
classical slice, and the shadow obstruction tower collapses to
class~$\mathbf{G}$ (Gaussian). Every non-trivial shadow class
($\mathbf{L}$, $\mathbf{C}$, $\mathbf{M}$) is invisible at
the PVA level. The PVA retains the $\lambda$-bracket (the
collision residues of $m_2$), which encodes the classical
$r$-matrix and hence the Drinfeld--Kohno data; the quartic and
higher shadow invariants ($S_4$, $Q^{\mathrm{contact}}$, the
infinite tower for class~$\mathbf{M}$) are quantum
corrections that require the chain-level $A_\infty$ structure.
\end{remark}

\begin{proof}[Completion of the proof of Theorem~\ref{thm:cohomology_PVA}]
Part~\textup{(i)} is Lemma~\ref{lem:lambda_descends}. Part~\textup{(ii)} is
Corollary~\ref{cor:PVA-structure}. These two parts are proved under
Definition~\ref{def:log-SC-algebra} together with
Theorem~\ref{thm:FM-calculus}; no higher-operation formality is used.
\end{proof}

\subsection{Two tangible studies}

\begin{example}[Free holomorphic--topological theory]
\label{ex:repaired-free-pva}
In the free theory there is no interaction vertex. Hence the binary kernel has no singular
part:
\[
m_2^{\mathrm{sing}}(a,b;\zeta)=0,
\]
and the regular part is the ordinary Wick product. Therefore
\[
\{a_\lambda b\}=0,
\qquad
[a]\cdot[b]=[\mu(a,b)]
\]
is simply the commutative product on the free operator algebra. Jacobi and Leibniz are
tautological, and the compatible contraction required by
Proposition~\ref{prop:m3_vanish} is immediate because all $m_{k\ge3}$ vanish already at
the chain level.
\end{example}

\begin{example}[Abelian current algebra]
\label{ex:repaired-abelian-current}
Let $J$ be the boundary current of the abelian Chern--Simons / free gauge-field example.
Its binary OPE kernel has singular part
\[
m_2^{\mathrm{sing}}(J,J;\zeta)=k\,\zeta^{-2},
\]
so the Borel transform gives the polynomial bracket
\[
\{J_\lambda J\}=k\,\lambda.
\]
The product is the regular constant term, traditionally written $J\cdot J= :JJ:$.
The proof above reproduces the standard PVA identities:
\begin{align*}
\{J_\lambda(J\cdot J)\}
&=
\{J_\lambda J\}\cdot J + J\cdot\{J_\lambda J\}
=2k\lambda\,J,\\
\{J_\lambda\{J_\mu J\}\}
&=0
=
\{\{J_\lambda J\}_{\lambda+\mu}J\}
+
\{J_\mu\{J_\lambda J\}\},
\end{align*}
because $\{J_\lambda J\}=k\lambda$ is central. This example makes the
dictionary completely concrete:
\[
k\,\zeta^{-2}\ \xleftrightarrow{\ \mathcal B\ }\ k\lambda.
\]
\end{example}

\begin{example}[Affine $\widehat{\mathfrak{sl}}_2$ at level~$k$; \ClaimStatusProvedHere]
\label{comp:pva-descent-affine-sl2}
The affine current algebra
$\widehat{\mathfrak{sl}}_2$ at level~$k\neq-2$ realises D2--D5 and the nullhomotopy D6 criterion explicitly.
The boundary chiral algebra~$\A$ is generated by currents
$e(z),f(z),h(z)$ with the standard affine Kac--Moody OPE
$J^a(z)J^b(w) \sim k\delta^{ab}/(z{-}w)^2 + f^{ab}{}_c\,J^c(w)/(z{-}w)$
at level~$k$ (see Kac~\cite{Kac98}, \S2.4, or
Frenkel--Ben-Zvi~\cite{FBZ04}, \S11.1).
In the Chevalley basis $(e,f,h)$, the singular OPE kernels are
\begin{align*}
m_2^{\mathrm{sing}}(e,f;\zeta)
 &= k\,\zeta^{-2}+h\,\zeta^{-1},
&
m_2^{\mathrm{sing}}(h,e;\zeta)
 &= 2e\,\zeta^{-1},
\\
m_2^{\mathrm{sing}}(h,f;\zeta)
 &= -2f\,\zeta^{-1},
&
m_2^{\mathrm{sing}}(h,h;\zeta)
 &= 2k\,\zeta^{-2}.
\end{align*}
All other generator pairings are regular.
Reading off the mode products
$e_{(0)}f=h$, $e_{(1)}f=k$, $h_{(0)}e=2e$, $h_{(0)}f=-2f$,
$h_{(1)}h=2k$
and applying the Borel transform
$\{a_\lambda b\}=\sum_{n\ge0}(a_{(n)}b)\,\lambda^n/n!$
gives the $\lambda$-bracket on $H^\bullet(\A,Q)$:
\begin{equation}
\label{eq:sl2-lambda-brackets}
\{e_\lambda f\}=h+k\lambda,
\qquad
\{h_\lambda e\}=2e,
\qquad
\{h_\lambda f\}=-2f,
\qquad
\{h_\lambda h\}=2k\lambda.
\end{equation}
The underlying commutative algebra is the differential polynomial ring
$R=\C[e^{(n)},f^{(n)},h^{(n)}:n\ge0]$
with $\partial(x^{(n)})=x^{(n+1)}$ and $x^{(0)}=x$.

\smallskip
\noindent\textbf{D2 (Sesquilinearity).}
We check both axioms on the pair $(h,e)$.
The mode relation $(\partial a)_{(n)} b=-n\,a_{(n-1)}b$ gives
\[
\{(\partial h)_\lambda e\}
=
\sum_{n\ge0}\frac{\lambda^n}{n!}\,(\partial h)_{(n)}e
=
\sum_{n\ge0}\frac{\lambda^n}{n!}\bigl(-n\,h_{(n-1)}e\bigr)
=
-\lambda\sum_{m\ge0}\frac{\lambda^m}{m!}\,h_{(m)}e
=
-\lambda\{h_\lambda e\}
=
-2\lambda\,e.
\]
Here the $n=0$ term vanishes and the re-indexing $m=n-1$ produces the
factor $-\lambda$. For the second axiom, the mode relation
$a_{(n)}(\partial b)=\partial(a_{(n)}b)+n\,a_{(n-1)}b$ yields
\begin{align*}
\{h_\lambda(\partial e)\}
&=
\sum_{n\ge0}\frac{\lambda^n}{n!}\bigl(\partial(h_{(n)}e)+n\,h_{(n-1)}e\bigr)
\\
&=
\partial\{h_\lambda e\}+\lambda\{h_\lambda e\}
=
(\partial+\lambda)\{h_\lambda e\}
=
(\partial+\lambda)(2e)
=
2\partial e+2\lambda\,e.
\end{align*}

\smallskip
\noindent\textbf{D3 (Skew-symmetry).}
The identity to establish is $\{f_\lambda e\}=-\{e_{-\lambda-\partial}f\}$.
Since all generators sit in degree~$0$,
Example~\ref{ex:repaired-abelian-current} already shows
that the skew-symmetry formula reduces to the standard PVA identity.
From the OPE $f(z)e(w)\sim k/(z{-}w)^2-h(w)/(z{-}w)$
we read off $f_{(0)}e=-h$, $f_{(1)}e=k$, so
\[
\{f_\lambda e\}=-h+k\lambda.
\]
On the other hand,
\[
-\{e_{-\lambda-\partial}f\}
=
-(h+k(-\lambda-\partial))
=
-h+k\lambda,
\]
because the standard PVA notation
$\{e_{-\lambda-\partial}f\}$ means that the substituted
$\partial$ is moved onto the coefficient of the corresponding
$\lambda$-power. Here that coefficient is the central scalar~$k$,
so $\partial k=0$. Thus no trailing $k\partial$ term survives, and
both sides equal $-h+k\lambda$.

\smallskip
\noindent\textbf{D4 (Jacobi identity for the $\lambda$-bracket).}
The PVA Jacobi identity
\begin{equation}
\label{eq:pva-jacobi-sl2-check}
\{a_\lambda\{b_\mu c\}\}
-
\{b_\mu\{a_\lambda c\}\}
=
\{\{a_\lambda b\}_{\lambda+\mu}c\}
\end{equation}
on the triple $(a,b,c)=(h,e,f)$ unfolds into three terms:
\begin{align*}
\{h_\lambda\{e_\mu f\}\}
&=
\{h_\lambda(h+k\mu)\}
=
\{h_\lambda h\}+0
=
2k\lambda,
\\[4pt]
\{e_\mu\{h_\lambda f\}\}
&=
\{e_\mu(-2f)\}
=
-2\{e_\mu f\}
=
-2(h+k\mu)
=
-2h-2k\mu,
\\[4pt]
\{\{h_\lambda e\}_{\lambda+\mu}f\}
&=
\{(2e)_{\lambda+\mu}f\}
=
2\{e_{\lambda+\mu}f\}
=
2\bigl(h+k(\lambda+\mu)\bigr)
=
2h+2k\lambda+2k\mu.
\end{align*}
The left-hand side of \eqref{eq:pva-jacobi-sl2-check} is
\[
2k\lambda-(-2h-2k\mu)
=
2h+2k\lambda+2k\mu
=
\text{RHS}.\qquad\checkmark
\]

As a second check, take $(a,b,c)=(e,f,e)$. Then:
\begin{align*}
\{e_\lambda\{f_\mu e\}\}
&=
\{e_\lambda(-h+k\mu)\}
=
-\{e_\lambda h\}
=
-(-2e)
=
2e,
\\[4pt]
\{f_\mu\{e_\lambda e\}\}
&=
\{f_\mu\,0\}
=
0,
\\[4pt]
\{\{e_\lambda f\}_{\lambda+\mu}e\}
&=
\{(h+k\lambda)_{\lambda+\mu}e\}
=
\{h_{\lambda+\mu}e\}
=
2e.
\end{align*}
The identity $2e-0=2e$ holds.\quad$\checkmark$

\smallskip
\noindent\textbf{D5 (Leibniz rule).}
The Leibniz identity reads $\{h_\lambda(e\cdot f)\}=\{h_\lambda e\}\cdot f+e\cdot\{h_\lambda f\}$,
where ``$\cdot$'' is the commutative product in the differential polynomial ring~$R$
(not the normally ordered product).
The right-hand side is
\[
\{h_\lambda e\}\cdot f+e\cdot\{h_\lambda f\}
=
2e\cdot f+e\cdot(-2f)
=
0.
\]
And $\{h_\lambda(e\cdot f)\}=0$ on the left because $\{h_\lambda-\}$ is a derivation
of~$R$ with $e\cdot f$ in weight~$0$ under the adjoint action
(the contributions from $e$ and $f$ cancel by opposite weights).
Both sides vanish.\quad$\checkmark$

For a nontrivial check, verify on the monomial $h^2\in R$:
\[
\{e_\lambda h^2\}
=
\{e_\lambda h\}\cdot h + h\cdot\{e_\lambda h\}
=
(-2e)\cdot h + h\cdot(-2e)
=
-4eh.
\]
Alternatively, by skew-symmetry and sesquilinearity,
$\{e_\lambda h\}=-\{h_{-\lambda-\partial}e\}=-2e$ (since $\{h_\mu e\}=2e$
is independent of~$\mu$, the substitution $\mu\mapsto-\lambda-\partial$
leaves it unchanged).
Thus $\{e_\lambda h^2\}=2\cdot(-2e)\cdot h = -4eh$.\quad$\checkmark$

\smallskip
\noindent\textbf{D6 (Vanishing under the affine nullhomotopy datum).}
The affine current algebra $\widehat{\mathfrak{sl}}_2$ arises as the
chiral algebra on the boundary of $3$d Chern--Simons theory, which
lies in the scope of Theorem~\ref{thm:physics-bridge} and therefore defines
a logarithmic $\SCchtop$-algebra with $d'=1$ topological direction.
On the chosen one-loop Koszul affine comparison model, the required
compatible topological nullhomotopies are part of the model: the
non-central simple-pole class is the binary affine factorisation
datum, not a new higher PVA operation. Under this additional datum,
Proposition~\ref{prop:m3_vanish} makes all $m_{k\ge3}$ $Q$-exact on
$Q$-closed inputs and hence zero in $H^\bullet(\A,Q)$. Concretely, the OPE of
$\widehat{\mathfrak{sl}}_2$ has poles of order at most~$2$, and the
affine nullhomotopy datum ensures that the corresponding $k$-point
integrals contribute only $Q$-boundaries.
\end{example}

\begin{remark}[Closure of the PVA chain and the D6 nullhomotopy hypothesis]
The arguments above close the foundational PVA chain and isolate the
D6 nullhomotopy criterion in a single geometric package:
\begin{itemize}
\item D2 is handled by the exchange cylinder and the unsymmetrized regular product;
\item D3 is handled by the same exchange cylinder, now in the singular sector and with
 the degree-$(-1)$ Koszul sign;
\item D4 is handled by the full three-face singular Stokes argument;
\item D5 is handled by the mixed regular-singular three-face Stokes argument;
\item D6 has the additional hypothesis of compatible topological
 nullhomotopies in addition to the factorisation
 $\widetilde\omega_k=\omega_k^{\mathrm{hol}}\wedge\omega_k^{\mathrm{top}}$
 from Definition~\ref{def:log-SC-algebra}.
\end{itemize}
D2--D5 and the vacuum axiom hold for any logarithmic
$\SCchtop$-algebra (Definition~\ref{def:log-SC-algebra}); the product
decomposition and FM calculus (Theorem~\ref{thm:FM-calculus}) supply
the required inputs. D6 holds on the additional nullhomotopic/formal
locus described above.
\end{remark}

\subsection{Functoriality of PVA descent}
\label{subsec:pva-functoriality}

The constructions above assemble into a functor.

\begin{definition}[Strict morphism of logarithmic $\SCchtop$-algebras]
\label{def:morphism-SC-algebra}
A \emph{strict morphism} $f \colon (A, Q_A, \partial_A,
\mathbf 1_A,\{m_k^A\})
\to (B, Q_B, \partial_B, \mathbf 1_B,\{m_k^B\})$ of logarithmic
$\SCchtop$-algebras is a $\hypAmbientPro$-continuous cochain map
$f \colon (A, Q_A) \to (B, Q_B)$
satisfying
\[
f\circ\partial_A=\partial_B\circ f,
\qquad
f(\mathbf 1_A)=\mathbf 1_B,
\]
and intertwining all logarithmic kernels coefficientwise:
\[
 f \circ m_k^A(a_1, \ldots, a_k; \lambda_1, \ldots, \lambda_{k-1})
 = m_k^B(f(a_1), \ldots, f(a_k); \lambda_1, \ldots, \lambda_{k-1})
\]
for all $k \geq 1$ and all $a_i \in A$.
Continuity means that $f$ commutes with the pro/logarithmic
completion, Laurent coefficient extraction, regular/singular
projection, and the polynomial Borel transform used in
Theorem~\textup{\ref{thm:cohomology-PVA-main}}.  Write
$\cat{SC\text{-}Alg}^{\log}$ for the category with these strict
unital morphisms.
Write $\cat{PVA}_{(-1)}$ for the category of
$(-1)$-shifted Poisson vertex algebras with PVA morphisms
(maps preserving $\mu$, $\{\cdot{}_\lambda\cdot\}$, and $\partial$).
\end{definition}

\begin{theorem}[Strict functoriality of PVA descent;
\ClaimStatusProvedHere; licensing $\gamma$ via the logarithmic
chain/pro ambient $\hypAmbientPro$ and strict unital coefficientwise
morphisms]
\label{thm:pva-descent-functor}
\index{PVA descent!functoriality|textbf}
The assignment
\[
 \mathcal{D} \colon \cat{SC\text{-}Alg}^{\log}
 \longrightarrow \cat{PVA}_{(-1)},
 \qquad
 A \longmapsto \bigl(H^\bullet(A,Q),\; \mu,\;
 \{\cdot{}_\lambda\cdot\},\; \partial\bigr),
\]
on objects defined by Theorem~\textup{\ref{thm:cohomology-PVA-main}},
extends to a functor: given a strict morphism
$f \colon A \to B$ of logarithmic $\SCchtop$-algebras,
$\mathcal{D}(f) := H^\bullet(f) \colon H^\bullet(A,Q_A) \to
H^\bullet(B,Q_B)$ is a morphism of $(-1)$-shifted PVAs.
\end{theorem}

\begin{proof}
The map $H^\bullet(f)$ preserves each piece of the PVA structure
because the morphism is strict, unital, and continuous in
$\hypAmbientPro$.

\medskip\noindent
\textbf{Step 1: Product.}
Since $f$ intertwines the binary logarithmic kernels and commutes with
regular projection and constant-term extraction,
\[
f\bigl(\mu_2^{A,\mathrm{reg}}(a,b;0)\bigr)
=\mu_2^{B,\mathrm{reg}}(f(a),f(b);0).
\]
As $f$ is a cochain map, it descends to
$H^\bullet(f)(\mu_A([a],[b])) = \mu_B(H^\bullet(f)[a], H^\bullet(f)[b])$.

\medskip\noindent
\textbf{Step 2: $\lambda$-bracket.}
The singular part $m_2^{\mathrm{sing}}(a,b;\zeta) = \sum_n
(a_{(n)}b)\,\zeta^{-n-1}$ is intertwined coefficientwise:
$f(a_{(n)}^A b) = f(a)_{(n)}^B f(b)$. Applying the Borel
transform, which is a linear operation on coefficients and commutes
with $f$ by $\hypAmbientPro$-continuity, gives
\[
 f\bigl(\{a {}_\lambda b\}_A\bigr)
 = f\Bigl(\sum_n \tfrac{\lambda^n}{n!}\, a_{(n)}^A b\Bigr)
 = \sum_n \tfrac{\lambda^n}{n!}\, f(a)_{(n)}^B f(b)
 = \{f(a) {}_\lambda f(b)\}_B.
\]
Descent to cohomology follows because $f$ is a cochain map, and
the $\lambda$-bracket descends to $H^\bullet$ by
Lemma~\ref{lem:lambda_descends}.

\medskip\noindent
\textbf{Step 3: Translation operator.}
$H^\bullet(f) \circ \partial_A = \partial_B \circ H^\bullet(f)$
on cohomology, since $f \circ \partial_A = \partial_B \circ f$
on cochains and $\partial$ commutes with $Q$.

\medskip\noindent
\textbf{Step 4: Unit.}
$f(\mathbf{1}_A) = \mathbf{1}_B$ by the strict unitality clause in
Definition~\textup{\ref{def:morphism-SC-algebra}}.

\medskip\noindent
\textbf{Step 5: Composition and identity.}
If $f \colon A \to B$ and $g \colon B \to C$ are morphisms,
then $H^\bullet(g \circ f) = H^\bullet(g) \circ H^\bullet(f)$
by functoriality of $H^\bullet$. Similarly
$H^\bullet(\mathrm{id}_A) = \mathrm{id}_{H^\bullet(A)}$.
\end{proof}

\begin{remark}[Quasi-isomorphism invariance]
\label{rem:pva-descent-qi}
If $f \colon A \xrightarrow{\sim} B$ is a quasi-isomorphism
of logarithmic $\SCchtop$-algebras, then $\mathcal{D}(f)$ is
an isomorphism of $(-1)$-shifted PVAs. In particular, the PVA
structure on $H^\bullet(A,Q)$ is a quasi-isomorphism invariant
of the $\SCchtop$-algebra, a necessary condition for it to
capture physically meaningful data.
\end{remark}

\section{K3 base-point PVA}
\label{sec:pvadr-supplement}

\begin{remark}[PVA descent at the K3 base point]
\label{rem:pvadr-1}
The K3 base-point PVA descent is the binary classical shadow of the
bi-based Ran factorisation. The bridge from the Siegel parameter base
$\bar{\mathcal A}_2$ to the chiral base
$\mathrm{Ran}(E^{\mathrm{nod,sm}}_{24})$ with $24$ Kodaira $I_1$
nodes is implemented by the averaging morphism
$\mathrm{av} = \mathrm{Torelli}_{K3}\circ\mathrm{KugaSatake}\circ j_{\mathrm{Kod}}$.
The PVA descent provides the BV classical limit
$\hbar_{\mathrm{BV}}\to 0$ of the modular-PVA quantisation at fixed
quantum-group fibre, with $R_{\mathrm{Sieg,dyn}}$ supplied by
Pasol--Zagier 2013
\cite{KugaSatake1967,PasolZagier2013}. Compare Vol.~III
Chapter~\ref{V3-ch:k3-chiral-bialgebra-platonic}.
\end{remark}

\begin{remark}[Conductor specialisation and the PVA limit]
\label{rem:pvadr-2}
The conductor identity fixes the quantum-group fibre, not the BV
classical parameter:
\[
\hbar_{\mathrm{QG}}^2\cdot K^{\kappaChHodge}=-1,\qquad
K^{\kappaChHodge}=8,\qquad
\hbar_{\mathrm{QG}}^2=-1/8.
\]
The PVA limit is instead the associated graded at
$\hbar_{\mathrm{BV}}=0$ of the BV deformation. Thus the K3
classical-limit PVA is formed at fixed Lusztig fibre
$\hbar_{\mathrm{QG}}^2=-1/8$ and then passed to
$\hbar_{\mathrm{BV}}\to0$; no identification of
$\hbar_{\mathrm{BV}}$ with the Lusztig value is asserted
\cite{Lusztig1990}.
Compare Chapter~\ref{V3-ch:k3-chiral-bialgebra-platonic}.
\end{remark}

\section{Weight-completed K3 PVA}
\label{sec:pvadr-supplement-ii}

\begin{remark}[Two-$\hbar$ convention bridge]
\label{rem:pvadr-two-hbar-bridge}
Two distinct $\hbar$'s appear in the K3 base-point PVA descent and should not
be conflated. (i) The \emph{BV--PVA} $\hbar_{\mathrm{BV}}$
(Remark~\ref{rem:pva-quasi-classical-bar-avatar}): classical limit
$\hbar_{\mathrm{BV}}\to 0$ extracts the binary operation $m_2$ and the
$(-1)$-shifted PVA bracket from $\Theta_{\mathcal A}$.
(ii)~The \emph{quantum-group} $\hbar_{\mathrm{QG}}$: the Etingof--Kazhdan
super-quantisation deformation parameter used to expand the twisted
Siegel--Borcherds associator
$\widetilde\Phi^{\mathrm{Sieg\text{-}Bor}}_{\mathrm{Sp}_4}$
as $1 + \sum_{n\ge 3}\hbar_{\mathrm{QG}}^n \phi^{(n)}$. The
Lusztig-specialisation locus $\hbar_{\mathrm{QG}}^2 = -1/8$
(Remark~\ref{rem:pvadr-2}) is a locus in the \emph{quantum-group}
parameter, and the classical-limit PVA is the
$\hbar_{\mathrm{BV}}\to 0$ shadow of the $A_\infty$-quasi-Hopf
structure at any fixed $\hbar_{\mathrm{QG}}$. The conductor identity
\[
 \hbar_{\mathrm{QG}}^{\,2}\,K^{\kappaChHodge}=-1,\qquad
 K^{\kappaChHodge}=8
\]
selects the quantum-group fibre. Identifying the two deformation
parameters would impose
$\hbar_{\mathrm{BV}}^2+1/8=0$, which is incompatible with the PVA
closed point $\hbar_{\mathrm{BV}}=0$. The comparison base is therefore
the product of the BV formal disk with the fixed Lusztig fibre, and
$\hbar_{\mathrm{BV}}$ remains the formal parameter of the
associated-graded PVA construction
\cite{EtingofKazhdan2000,Lusztig1990}.
\end{remark}

\begin{remark}[Classical-limit PVA on finite Hall windows]
\label{rem:pvadr-classical-limit}
For each $N$, let $\mathbf H_{\Delta_5}^{\le N}$ be a finite
recognized Hall window in the sense of
Definition~\ref{def:k3-finite-hall-windows}.  Let
$\mathcal A_N$ be the logarithmic $\SCchtop$-algebra obtained from
that finite Hall double after the radical quotient
$\operatorname{Rad}_N$ and the source-recognition map $\rho_N$ have
been fixed.  The $\hbar_{\mathrm{BV}}\to0$ classical limit is formed
first on this finite object:
\[
 \mathcal D_N
 :=
 \bigl(
 H^\bullet(\mathcal A_N,Q_N),\,
 \mu_N,\,
 \{-_\lambda-\}_N
 \bigr).
\]
Theorem~\ref{thm:cohomology-PVA-main} applies to $\mathcal A_N$ and
gives a $(-1)$-shifted PVA $\mathcal D_N$.  No PVA on the completed
object is asserted before these finite descents and their transition
maps are supplied.
\end{remark}

\begin{proposition}[Windowwise PVA descent and recognized pro-limit;
\ClaimStatusConditional; licensing $\alpha+\beta+\gamma$ via the
height-$N$ Hall--Borcherds recognition datum, the maps $\rho_N$,
strict $\hypAmbientPro$ morphisms, and the weight-completed ambient
$\hypAmbientWtCpl$]
\label{rem:pvadr-MC-integrability}
\label{prop:pvadr-windowwise-descent}
Assume that the height-$N$ Hall--Borcherds recognition package of
Definition~\ref{def:k3-finite-hall-windows} is fixed.  Then
$\mathcal D_N$ is a $(-1)$-shifted PVA.  If the transition maps
$\mathbf H_{\Delta_5}^{\le N+1}\to\mathbf H_{\Delta_5}^{\le N}$
are strict $\hypAmbientPro$-continuous morphisms of logarithmic
$\SCchtop$-algebras in the sense of
Definition~\ref{def:morphism-SC-algebra}, carry
$\operatorname{Rad}_{N+1}$ into $\operatorname{Rad}_N$, and commute
with the recognition maps $\rho_{N+1}$ and $\rho_N$, then in
$\hypAmbientWtCpl$
\[
 \widehat{\mathcal D}_{\Delta_5}
 :=
 \widehat{\varprojlim}_{N}\,\mathcal D_N
\]
is a completed $(-1)$-shifted PVA.  Without those compatible finite
recognition maps, the pro-limit PVA is not part of the theorem.
\end{proposition}

\begin{proof}
The finite Hall double has associative Hall multiplication,
coassociative Hall coproduct, opposite half, Cartan part, and finite
Hall pairing by the recognition package.  The radical
$\operatorname{Rad}_N$ is included as a Hopf ideal in
Definition~\ref{def:k3-finite-hall-windows}; in the
$\hbar_{\mathrm{BV}}\to0$ associated graded it is therefore a
Poisson/PVA ideal.  The radical quotient $\mathcal A_N$ inherits the
binary regular product and the singular residue bracket from the finite
Hall double.  The cohomological descent theorem
(Theorem~\ref{thm:cohomology-PVA-main}) gives the PVA identities on
$H^\bullet(\mathcal A_N,Q_N)$.

Jacobi, Leibniz, and skew-symmetry are inherited from the finite Hall
double and then from the quotient; no Kostant--Kumar resolution or
Kac classification theorem is used for the BKM super-Manin pair.  The
imaginary-root contribution enters only through the finite recognition
map $\rho_N$, which compares the windowed Hall constants with the
$W^{\le N}$ Borcherds coefficients of the Gritsenko--Nikulin
denominator.  The strict transition maps are sent by
Theorem~\ref{thm:pva-descent-functor} to PVA morphisms
$\mathcal D_{N+1}\to\mathcal D_N$.  The completed inverse limit is
then computed relationwise in $\hypAmbientWtCpl$: every PVA identity
in the limit is the compatible inverse system of the corresponding
finite-window identities.
\end{proof}

\begin{remark}[Completed K3 PVA as conditional inverse limit]
\label{rem:pvadr-weight-completed}
Under the hypotheses of Proposition~\ref{prop:pvadr-windowwise-descent}
the recognized inverse-limit Lie bialgebra may be written
\[
 \widehat{\mathfrak g}^{\mathrm{super},\mathrm{rec}}_{\Delta_5}
 :=
 \widehat{\varprojlim}_{N}\,
 \left(
 \mathfrak g_{\Delta_5}^{\mathrm{Hall},\le N}
 \right)_{\!/\operatorname{Rad}_N}.
\]
This is the completed source-recognized object whose PVA is
$\widehat{\mathcal D}_{\Delta_5}$.  The Gritsenko lift of
$\eta^9\theta_1$ to $\Delta_5$ supplies the automorphic denominator
target and the finite Borcherds coefficient windows; it does not by
itself supply the Hall product, coproduct, radical quotient, or
transition maps.  On the strict non-completed ambient the infinite
imaginary-root sector still prevents a global PVA statement.  On the
completed ambient the statement is conditional on the windowwise
source recognition just stated.  Compare Vol.~I
Remark~\ref{V1-rem:nc-Ainfty-open} and Vol.~III
Chapter~\ref{V3-ch:k3-chiral-bialgebra-platonic}.
\cite{GritsenkoNikulin1998,Gritsenko1999,Lusztig1990,KhanZeng2025}.
\end{remark}
